\documentclass[11pt]{amsart}

\usepackage[T1]{fontenc}
\usepackage{lmodern}
\usepackage{amsmath,amssymb,mathtools}
\usepackage{microtype}
\usepackage{xcolor}
\usepackage[pdfusetitle,colorlinks=true,linkcolor=blue!45!black,citecolor=blue!45!black,urlcolor=blue!45!black]{hyperref}
\raggedbottom

\newtheorem{theorem}{Theorem}[section]
\newtheorem{proposition}[theorem]{Proposition}
\newtheorem{lemma}[theorem]{Lemma}
\newtheorem{corollary}[theorem]{Corollary}
\theoremstyle{definition}
\newtheorem{definition}[theorem]{Definition}
\theoremstyle{remark}
\newtheorem{remark}[theorem]{Remark}

\DeclareMathOperator{\Hess}{Hess}
\DeclareMathOperator{\Isom}{Isom}
\DeclareMathOperator{\rank}{rank}
\DeclareMathOperator{\tr}{tr}
\DeclareMathOperator{\spanop}{span}
\DeclareMathOperator{\secop}{sec}

\newcommand{\R}{\mathbb R}
\newcommand{\Sph}{\mathbb S}
\newcommand{\ip}[2]{\left\langle #1,#2\right\rangle}
\newcommand{\norm}[1]{\left\lVert #1\right\rVert}
\newcommand{\restr}[2]{\left.#1\right|_{#2}}

\title[Curvature obstructions and flat tori]{Curvature Obstructions and Flat Totally Geodesic Tori in Hessian Deformations of Products of Spheres}
\author{Lars R\"onnb\"ack}
\address{Department of Computer and Systems Sciences, Stockholm University, Stockholm, Sweden}
\email{lars@uptochange.com}
\date{}

\subjclass[2020]{53C20, 53C21, 53C24}
\keywords{mixed sectional curvature, product of spheres, Hopf conjecture, Hessian deformation, metric two-jet, curvature threshold, totally geodesic flat torus}

\begin{document}

\begin{abstract}
Let $h$ be the product of the unit round metrics on
$\Sph^m\times\Sph^n$, let $A:\R^{n+1}\to\R^{m+1}$ be linear, and put
\[
 s_A(x,y)=\ip{x}{Ay},\qquad
 g_{\gamma,A}=h+\gamma\Hess_h s_A.
\]
Writing $\sigma_1\geq\sigma_2\geq\cdots\geq0$ for the singular values
of $A$, padded by zeros, we prove that $g_{\gamma,A}$ is Riemannian
exactly when $|\gamma|(\sigma_1+\sigma_2)<1$.  If $m,n\geq2$ and
$A\neq0$, the top three singular pairs determine a totally geodesic
$\Sph^2\times\Sph^2$ containing a plane of negative sectional
curvature at every nonzero Riemannian parameter.  Thus the round
product is the unique nonnegatively curved metric on each such
Hessian ray.  We also obtain an exact formula for every
background-mixed curvature and an explicit dimension-free negative
margin on normalized parameter annuli.  For every pair of singular
directions, the metric contains a flat totally geodesic $2$-torus, so
no member of the family has strictly positive sectional curvature.

For arbitrary paths $h+t\Hess f+t^2B+O(t^3)$, we give a two-jet normal
form modulo diffeomorphisms.  An integrated trace identity forces every
quadratic mixed repair inequality to be an equality on a closed
product.  We classify the resulting coupled kernel modulo gauge and
factor variations: one tensor leg must be Killing and the other
co-closed.  On $\Sph^2\times\Sph^2$, every nonzero coupled class has a
nearby plane with curvature $-\delta s^2+O(s^3)$; this includes
cross-degree superpositions, simultaneous mixtures of both one-sided
towers, and the full $(1,1)$ block.  The associated canonical
corrected Hessian path therefore fails at order $t^4$.  A transverse
rotational sum also admits an exact negative oblique formula.  These
are obstructions for the stated Hessian family and kernel rays, not a
proof of the Hopf conjecture.
\end{abstract}

\maketitle

\section{Introduction}

The question whether $\Sph^2\times\Sph^2$ admits a metric of strictly
positive sectional curvature is the classical Hopf conjecture.  It is
open.  A useful way to approach the problem without overstating what a
special construction can prove is to choose a geometrically natural
family of metrics, determine its full parameter range, and identify an
exact obstruction inside that family.

This paper carries out that program for the finite-dimensional family
\begin{equation}\label{eq:family-intro}
 g_{\gamma,A}=h+\gamma\Hess_h s_A,
 \qquad s_A(x,y)=\ip{x}{Ay},
\end{equation}
on a product of unit spheres.  Here
$A:\R^{n+1}\to\R^{m+1}$ is arbitrary and $\gamma\in\R$.  The functions
$s_A$ are precisely the linear combinations of bilinear products of
degree-one spherical harmonics.  Thus the family includes couplings with
no continuous symmetry inherited from the two round factors, but it
retains enough structure for the exact global results below.

The results have three parts.  First, a sharp positivity criterion
depends on the two largest singular values of $A$.  Second, an exact
mixed-curvature formula gives a spectral-gap obstruction: when
$\sigma_1>\sigma_2$, a fixed mixed plane has negative curvature at
every nonzero point of the Riemannian ray.  A top-critical oblique
plane in the totally geodesic top-three singular
$\Sph^2\times\Sph^2$ is negative for every matrix of rank at least
two; a zero-critical oblique plane handles rank one and lies in the
same slice.  Thus all dimensions are detected by a four-dimensional
submodel, with an explicit dimension-free quantitative margin.
Third, each pair of
singular directions cuts out a product of great circles which remains
totally geodesic and becomes intrinsically flat for the deformed
metric.  Together the two curvature mechanisms show that, whenever
$m,n\geq2$ and $A\neq0$, the round product is the only metric on its
Hessian ray with nonnegative sectional curvature.  At the isotropic
$\Sph^2\times\Sph^2$ endpoint all background-mixed
curvatures are nevertheless nonnegative.  This mixed--oblique
dichotomy shows in particular that background-mixed nonnegativity is
not a proxy for full nonnegative sectional curvature.

The exact rays also admit a useful deformation-theoretic
interpretation.  Since
$\Hess f=\frac12\mathcal L_{\nabla f}h$, their first variation is
pure gauge.  We give a two-jet normal form for arbitrary corrected
paths $h+t\Hess f+t^2B+O(t^3)$ and an explicit mixed-curvature
functional which every nonnegatively curved path must make
nonnegative.  Its integrated trace vanishes, so on the closed product
all those values must in fact be zero.  For $f=s_A$, singular triples
turn this into concrete identities for the second-order correction
$B$.  The canonical gauge correction satisfies them, whereas the
uncorrected ray $B=0$ violates one for every $A\neq0$.
Modulo diffeomorphisms and factor variations, we classify the full
kernel: every coupled class is represented by a pure mixed,
divergence-free tensor with a Killing one-form on at least one factor
and a co-closed one-form on the other.  On $\Sph^2\times\Sph^2$ this
leaves exactly the harmonic bidegrees $(1,\ell)$ and $(\ell,1)$.  A
two-curl moving-plane calculation and a frame-average spectral gap
obstruct every nonzero element of this entire coupled quotient,
including arbitrary cross-degree and two-tower mixtures.  A signed
singular-value argument closes the full $(1,1)$ block.  The canonical
transverse rotational sum additionally has an exact negative oblique
formula.

\begin{theorem}\label{thm:main}
Let $m,n\geq1$, set
$q=\min\{m+1,n+1\}$, and let
$\sigma_1\geq\sigma_2\geq\cdots\geq\sigma_q\geq0$ be the singular
values of $A:\R^{n+1}\to\R^{m+1}$.  If $A\neq0$, then
\[
 g_{\gamma,A}\text{ is Riemannian}
 \quad\Longleftrightarrow\quad
 |\gamma|(\sigma_1+\sigma_2)<1.
\]
For every $\gamma$ in this interval and every choice of singular
frames, $g_{\gamma,A}$ contains one flat totally geodesic torus for
each pair $1\leq i<j\leq q$.  In particular,
\[
 \secop_{g_{\gamma,A}}(T_pT_{ij})=0
 \quad\text{for every }p\in T_{ij},
\]
and $g_{\gamma,A}$ never has strictly positive sectional curvature.
If moreover $m,n\geq2$, then the totally geodesic
$\Sph^2\times\Sph^2$ determined by the top three singular pairs
contains a plane of strictly negative sectional curvature for every
nonzero Riemannian parameter.  Hence the metric does not have
nonnegative sectional curvature.  If $\sigma_2>0$, the plane can be
chosen top-critical oblique; in rank one it can be chosen
zero-critical oblique.  If $\sigma_1>\sigma_2$, it can also be chosen
background-mixed.
If $A=0$, then $g_{\gamma,A}=h$ for every $\gamma$.
\end{theorem}

The negative mixed curvature in the examples below is strict.  It
therefore persists under small changes of the matrix, base point, and
$2$-plane; it is not an isolated algebraic coincidence.

For $\Sph^2\times\Sph^2$, a chosen singular frame gives three such
tori.  When $A$ is invertible with three distinct singular values, the
factor-preserving isometries inherited from the round product form a
finite group.  This matters because the powerful theorem of
Hsiang--Kleiner excludes positive curvature on
$\Sph^2\times\Sph^2$ in the presence of an isometric circle action
\cite{HsiangKleiner1989}.  Our obstruction is instead visible directly
on fixed tori of finite involutions.  We do not claim that the full
isometry group of the deformed metric is always finite.

Deformations of product metrics have long been studied in connection
with Hopf's question, notably by Bourguignon
\cite{Bourguignon1975,BourguignonDeschampsSentenac1973}.  More recent
work has produced obstructions under symmetry assumptions
\cite{HsiangKleiner1989,AmannKennard2017} and within conformal or warped
classes \cite{SeckNiangMasiello2022}; see also the survey
\cite{Amann2021}.  The present result is deliberately narrower and
fully explicit: it gives a complete positive-curvature obstruction
for the family~\eqref{eq:family-intro}, not the Hopf conjecture.  Its
additional contribution is that the same
mechanism works on every product $\Sph^m\times\Sph^n$ and identifies
the exact Riemannian interval rather than only a perturbative
neighborhood of the product metric.

The proof is organized as follows.  Section~\ref{sec:hessian} gives an
intrinsic formula for the Hessian and its tangent-space block form.
Section~\ref{sec:positive} proves the sharp positivity interval,
derives the exact mixed-curvature formula, proves the spectral-gap and
top-critical oblique obstructions, isolates the
isotropic mixed--oblique dichotomy, and analyzes plane-wise
sign-change thresholds.
Section~\ref{sec:tori} constructs and analyzes the flat totally
geodesic tori.  Section~\ref{sec:symmetry} separates the inherited
factor-preserving symmetry from the generally harder full isometry
group.  Section~\ref{sec:two-jet} gives the two-jet normal form and
the integrated rigidity, explicit repair identities, the full coupled
kernel classification, the nonlinear obstruction for its entire
coupled quotient, and an exact transverse rotational model.
Section~\ref{sec:scope} records
the precise relation to the Hopf problem and to possible extensions.

\section{The Hessian deformation}\label{sec:hessian}

Let
\[
 M=\Sph^m\times\Sph^n
 \subset\R^{m+1}\times\R^{n+1},
 \qquad h=h_m\oplus h_n,
\]
where both factors have radius one.  At $(x,y)\in M$ we identify
\[
 T_{(x,y)}M=x^\perp\oplus y^\perp.
\]
The Levi-Civita connection and Hessian below are those of $h$.

\begin{proposition}[Hessian formula]\label{prop:hessian}
For $X,X'\in x^\perp$ and $Y,Y'\in y^\perp$,
\begin{align}
 \Hess s_A\big((X,Y),(X',Y')\big)
 ={}&-s_A(x,y)\big(\ip{X}{X'}+\ip{Y}{Y'}\big) \notag\\
 &+\ip{X}{AY'}+\ip{X'}{AY}.       \label{eq:hessian}
\end{align}
Equivalently, with $P_x$ and $P_y$ the orthogonal projections onto
$x^\perp$ and $y^\perp$, respectively,
\begin{equation}\label{eq:block}
 g_{\gamma,A}=
 \begin{pmatrix}
  \lambda I_{x^\perp} & \gamma P_xAP_y\\
  \gamma P_yA^{\mathsf T}P_x & \lambda I_{y^\perp}
 \end{pmatrix},
 \qquad \lambda=1-\gamma s_A(x,y).
\end{equation}
\end{proposition}

\begin{proof}
For a fixed vector $a\in\R^{m+1}$, the restriction of the linear
function $x\mapsto\ip{x}{a}$ to the unit sphere satisfies
\begin{equation}\label{eq:linear-hess}
 \Hess_{\Sph^m}\ip{x}{a}=-\ip{x}{a}\,h_m.
\end{equation}
Indeed, the tangential gradient is $a-\ip{x}{a}x$, and differentiating
it with the Gauss formula gives~\eqref{eq:linear-hess}.  Holding $y$
fixed and taking $a=Ay$ yields the first pure block of
\eqref{eq:hessian}; holding $x$ fixed gives the second.  Since the
product connection differentiates independently in the two factors,
the mixed derivative is
\[
 X\big(\ip{x}{AY'}\big)=\ip{X}{AY'}.
\]
Symmetry of the Hessian supplies the other mixed entry.  Adding $h$
gives~\eqref{eq:block}.
\end{proof}

Formula~\eqref{eq:block} already exhibits the two quantities that
compete for definiteness: the scalar diagonal coefficient
$1-\gamma s_A$ and the largest singular value of the tangent
compression $P_xAP_y$.  The next section optimizes their combination
globally.

\section{The sharp Riemannian range and mixed curvature}
\label{sec:positive}

We use a two-frame form of the Ky Fan variational principle.  Singular
values are always listed in nonincreasing order and padded by zeros.

\begin{lemma}[Two-frame principle]\label{lem:kyfan}
Let $A:\R^r\to\R^p$ have singular values
$\sigma_1\geq\sigma_2\geq\cdots$.  Then
\begin{equation}\label{eq:kyfan}
 \sup_{\substack{x,X\in\R^p\text{ orthonormal}\\
                  y,Y\in\R^r\text{ orthonormal}}}
 \left(\ip{x}{Ay}+\left|\ip{X}{AY}\right|\right)
 =\sigma_1+\sigma_2.
\end{equation}
The same supremum is obtained with $-\ip{x}{Ay}$ in place of
$\ip{x}{Ay}$.
\end{lemma}

\begin{proof}
Change the sign of $X$ if necessary so that the second pairing is
nonnegative.  The expression is then the trace pairing of $A$ with a
rank-two partial isometry.  The Ky Fan trace inequality bounds it by
$\sigma_1+\sigma_2$; equivalently, this is the rank-two case of the
variational characterization of the nuclear norm of a compression
\cite[Chapter~III]{Bhatia1997}.
Equality is obtained from left and right singular vectors
$x=u_1$, $X=u_2$, $y=v_1$, and $Y=v_2$.  Replacing $x$ by $-x$
proves the final assertion.
\end{proof}

For completeness, define the largest tangent coupling at $(x,y)$ by
\begin{equation}\label{eq:tau}
 \tau(x,y)=
 \max_{\substack{X\perp x,\;Y\perp y\\
                  \norm{X}=\norm{Y}=1}}
 \left|\ip{X}{AY}\right|.
\end{equation}
Lemma~\ref{lem:kyfan} says precisely that
\begin{equation}\label{eq:two-sups}
 \sup_{x,y}\big(s_A(x,y)+\tau(x,y)\big)
 =\sup_{x,y}\big(\tau(x,y)-s_A(x,y)\big)
 =\sigma_1+\sigma_2.
\end{equation}

\begin{proposition}[Exact definiteness interval]\label{prop:positive}
If $A\neq0$, then $g_{\gamma,A}$ is positive definite at every point
of $M$ if and only if
\begin{equation}\label{eq:pd-range}
 |\gamma|(\sigma_1+\sigma_2)<1.
\end{equation}
At either endpoint the tensor is degenerate somewhere, and beyond an
endpoint it is not positive semidefinite everywhere.
\end{proposition}

\begin{proof}
At a fixed $(x,y)$, choose singular frames for the tangent compression
$P_xAP_y$.  The block matrix~\eqref{eq:block} is positive definite if
and only if
\begin{equation}\label{eq:pointwise-pd}
 \lambda>|\gamma|\tau(x,y),
 \qquad \lambda=1-\gamma s_A(x,y).
\end{equation}
The strict inequality also implies $\lambda>0$.

Suppose first that $\gamma>0$.  By~\eqref{eq:two-sups},
\[
 \lambda-\gamma\tau
 =1-\gamma(s_A+\tau)
 \geq1-\gamma(\sigma_1+\sigma_2).
\]
For $\gamma<0$, writing $a=|\gamma|$ gives
\[
 \lambda-a\tau
 =1-a(\tau-s_A)
 \geq1-a(\sigma_1+\sigma_2).
\]
Thus~\eqref{eq:pd-range} implies positive definiteness.

Conversely, take singular vectors satisfying
$Av_k=\sigma_k u_k$.  For $\gamma>0$, at $(x,y)=(u_1,v_1)$ evaluate
the metric on $(u_2,-v_2)$.  Proposition~\ref{prop:hessian} gives
\begin{equation}\label{eq:failure-plus}
 g_{\gamma,A}\big((u_2,-v_2),(u_2,-v_2)\big)
 =2\big(1-\gamma(\sigma_1+\sigma_2)\big).
\end{equation}
For $\gamma=-a<0$, use $(x,y)=(-u_1,v_1)$ and the vector
$(u_2,v_2)$ to obtain
\begin{equation}\label{eq:failure-minus}
 g_{-a,A}\big((u_2,v_2),(u_2,v_2)\big)
 =2\big(1-a(\sigma_1+\sigma_2)\big).
\end{equation}
These expressions vanish at the endpoints and are negative beyond
them.  Padding the singular values by zero makes the argument include
rank-one maps without modification.
\end{proof}

\begin{remark}
The occurrence of $\sigma_1+\sigma_2$, rather than only the operator
norm $\sigma_1$, is geometric.  One top singular direction fixes the
base point $(x,y)$ and contributes through $s_A$; a second orthogonal
direction remains available in the tangent spaces and contributes
through the mixed Hessian block.
\end{remark}

We next calculate every plane spanned by a vector from each background
factor.  The formula also explains why a nonnegative second-order term
does not imply nonnegative curvature for small nonzero parameters.

\begin{theorem}[Exact background-mixed curvature]
\label{thm:mixed-curvature}
Fix $p=(x,y)$ and put
\[
 a=P_xAy\in x^\perp,\qquad b=P_yA^{\mathsf T}x\in y^\perp.
\]
For $X\in x^\perp$ and $Y\in y^\perp$, set
\[
 r=\norm{X}^2,\quad w=\norm{Y}^2,\quad
 \alpha=\ip{a}{X},\quad \beta=\ip{b}{Y},
\]
and, using $h$ to identify vectors and covectors, define
\begin{align}
 \eta_{XX}&=\big(\tfrac r2a-\alpha X,-\tfrac r2b\big),\notag\\
 \eta_{XY}&=\big(-\tfrac\beta2X,-\tfrac\alpha2Y\big),
                                                        \label{eq:etas}\\
 \eta_{YY}&=\big(-\tfrac w2a,\tfrac w2b-\beta Y\big).\notag
\end{align}
Whenever $g_{\gamma,A}$ is Riemannian,
\begin{equation}\label{eq:exact-mixed-curvature}
 \secop_{g_{\gamma,A}}(X\wedge Y)
 =\frac{\gamma^2}{\Delta_{\gamma,A}(X,Y)}
 \left(
  g_{\gamma,A}^{-1}(\eta_{XY},\eta_{XY})
  -g_{\gamma,A}^{-1}(\eta_{XX},\eta_{YY})
 \right),
\end{equation}
where
\[
 \Delta_{\gamma,A}(X,Y)
 =g_{\gamma,A}(X,X)g_{\gamma,A}(Y,Y)
  -g_{\gamma,A}(X,Y)^2.
\]
In particular, if $X$ and $Y$ are $h$-unit vectors, then
\begin{equation}\label{eq:mixed-expansion}
 \secop_{g_{\gamma,A}}(X\wedge Y)
 =\frac{\gamma^2}{4}
 \left(\norm{a}^2-\alpha^2+\norm{b}^2-\beta^2\right)
 +O(\gamma^3).
\end{equation}
\end{theorem}

\begin{proof}
Write $H=\Hess_hs_A$ and choose $h$-normal coordinates at $p$ adapted
to $x^\perp\oplus y^\perp$.  If $\Gamma^g_{WUV}$ denotes the
Christoffel symbol of the first kind, then at $p$
\[
 \begin{aligned}
  \Gamma^g_{WUV}&=\gamma\mathcal C(U,V,W),\\
  \mathcal C(U,V,W)
  &=\frac12\big((\nabla_UH)(V,W)+(\nabla_VH)(U,W)\\
  &\hspace{42mm}-(\nabla_WH)(U,V)\big).
 \end{aligned}
\]
In the lower-index coordinate formula for curvature, the part involving
second derivatives of the metric is affine in $\gamma$.  Its constant
term vanishes on a background-mixed plane because the round product has
zero mixed curvature.  Its linear term vanishes as well: indeed,
$H=\frac12\mathcal L_{\nabla s_A}h$, so the linearized curvature is a
Lie derivative of the product curvature tensor.  In each term evaluated
on $(X,Y,X,Y)$ at least one vector remains in each factor, and the
product curvature tensor vanishes.  Consequently the exact lower-index
curvature is only the Christoffel-quadratic part,
\begin{equation}\label{eq:mixed-lower}
 R^g(X,Y,X,Y)=\gamma^2\left(
  g^{-1}(\mathcal C_{XY},\mathcal C_{XY})
  -g^{-1}(\mathcal C_{XX},\mathcal C_{YY})\right).
\end{equation}

It remains to calculate the three covectors.  Differentiating
Proposition~\ref{prop:hessian} with $U=(U_1,U_2)$ gives
\begin{align*}
 (\nabla_UH)(V,W)={}&-ds_A(U)h(V,W)
 -\ip{U_1}{V_1}\ip{b}{W_2}
 -\ip{U_2}{W_2}\ip{a}{V_1}\\
 &-\ip{U_1}{W_1}\ip{b}{V_2}
 -\ip{U_2}{V_2}\ip{a}{W_1}.
\end{align*}
Substitution of $(X,X)$, $(X,Y)$, and $(Y,Y)$ yields exactly the
covectors in~\eqref{eq:etas}.  Dividing~\eqref{eq:mixed-lower} by the
squared $g$-area proves~\eqref{eq:exact-mixed-curvature}.  Setting
$g^{-1}=h^{-1}+O(\gamma)$ and using $r=w=1$ gives
\eqref{eq:mixed-expansion}.
\end{proof}

The nonnegative coefficient in~\eqref{eq:mixed-expansion} does not
control the sign on a punctured neighborhood of the product metric.
In fact, a simple largest singular value forces strict negative
curvature along the entire punctured Riemannian ray.

\begin{theorem}[Generic spectral-gap obstruction]
\label{thm:spectral-gap}
Assume $m,n\geq2$ and $\sigma_1>\sigma_2$.  Write
\[
 s=\sigma_1,\qquad t=\sigma_2,\qquad \rho=\sigma_3,
\]
where singular values are padded by zeros, and set
\begin{equation}\label{eq:spectral-cd}
 c^2=\frac{s^2+t^2-2\rho^2}{2(s^2-\rho^2)},
 \qquad
 d^2=\frac{s^2-t^2}{2(s^2-\rho^2)}.
\end{equation}
Choose positive square roots and singular pairs satisfying
$Av_i=\sigma_i u_i$ and $A^{\mathsf T}u_i=\sigma_i v_i$, and define
\begin{equation}\label{eq:spectral-plane}
 \begin{aligned}
  x&=cu_1+du_3,& y&=v_2,& \mathsf X&=(u_2,0),\\
  B^2&=\frac{s^2+t^2}{2},&
  Y&=\frac{scv_1+\rho d v_3}{B},&
  \mathsf Y&=(0,Y).
 \end{aligned}
\end{equation}
Then $\mathsf X,\mathsf Y$ are $h$-orthonormal at $p=(x,y)$ and,
for every $\gamma$ in the Riemannian interval,
\begin{equation}\label{eq:spectral-gap-curvature}
 \secop_{g_{\gamma,A}}(\mathsf X\wedge\mathsf Y)
 =-\frac{
  \gamma^4(s^2-t^2)^2(s^2+t^2-2\rho^2)
 }{
  16(s^2+t^2)
  \left(1-\frac{\gamma^2}{2}(s^2-t^2+2\rho^2)\right)
 }.
\end{equation}
In particular, this curvature is strictly negative for every
$\gamma\neq0$ in that interval.
\end{theorem}

\begin{proof}
The numbers in~\eqref{eq:spectral-cd} are positive and add to one.
At the point and plane in~\eqref{eq:spectral-plane}, the quantities in
Theorem~\ref{thm:mixed-curvature} are
\[
 a=tu_2,\qquad b=BY,\qquad \alpha=t,\qquad \beta=B.
\]
Put
\[
 Z=du_1-cu_3,\qquad
 k=sdv_1-\rho cv_3,
\]
and abbreviate
\begin{equation}\label{eq:spectral-mu-kappa}
 \mu=\ip{k}{Y}=\frac{cd(s^2-\rho^2)}{B},
 \qquad
 \kappa^2=\norm{k}^2
 =\frac{s^2-t^2+2\rho^2}{2}.
\end{equation}
Since $s_A(x,y)=0$, the diagonal coefficient of the metric at $p$ is
one.  The tangent compression $C=P_xAP_y$ satisfies
\[
 C^{\mathsf T}u_2=0,
 \qquad
 Cw=\ip{k}{w}Z
 \quad\text{for }w\in\spanop\{v_1,v_3\}.
\]
The remaining singular directions form orthogonal blocks and do not
couple to $u_2$ or $Y$.  The Schur-complement formula, or a direct
inversion of the block spanned by $Z$ and $k$, therefore gives
\begin{align}
 g_{\gamma,A}^{-1}((u_2,0),(u_2,0))&=1,\notag\\
 g_{\gamma,A}^{-1}((0,Y),(0,Y))
 &=1+\frac{\gamma^2\mu^2}{1-\gamma^2\kappa^2}
 =:q_\gamma,                                      \label{eq:spectral-inverse}\\
 g_{\gamma,A}^{-1}((u_2,0),(0,Y))&=0.\notag
\end{align}
The denominator in~\eqref{eq:spectral-inverse} is positive throughout
the Riemannian interval; this also follows directly from positive
definiteness of the corresponding principal block.

For this plane the three covectors in~\eqref{eq:etas} reduce to
\[
 \eta_{XX}=\eta_{YY}
 =\left(-\frac t2u_2,-\frac B2Y\right),
 \qquad
 \eta_{XY}
 =\left(-\frac B2u_2,-\frac t2Y\right).
\]
Moreover, the squared $g_{\gamma,A}$-area of the plane is one.  Hence
Theorem~\ref{thm:mixed-curvature} yields
\begin{align*}
 \secop_{g_{\gamma,A}}(\mathsf X\wedge\mathsf Y)
 &=\frac{\gamma^2}{4}
   \left(B^2+t^2q_\gamma-t^2-B^2q_\gamma\right)\\
 &=-\frac{\gamma^4(B^2-t^2)\mu^2}
          {4(1-\gamma^2\kappa^2)}.
\end{align*}
Finally,
\[
 B^2-t^2=\frac{s^2-t^2}{2},
 \qquad
 \mu^2=
 \frac{(s^2+t^2-2\rho^2)(s^2-t^2)}{2(s^2+t^2)}.
\]
Substitution gives~\eqref{eq:spectral-gap-curvature}.  Every factor in
its numerator is positive when $s>t\geq\rho$, which proves the sign
claim.
\end{proof}

\begin{corollary}[Generic isolation from nonnegative curvature]
\label{cor:generic-isolation}
For $m,n\geq2$, the matrices with $\sigma_1>\sigma_2$ form an open
dense subset of
$\operatorname{Hom}(\R^{n+1},\R^{m+1})$.  For every such $A$,
\[
 \left\{\gamma:
  g_{\gamma,A}\text{ is Riemannian and }\secop_{g_{\gamma,A}}\geq0
 \right\}=\{0\}.
\]
Thus the round product is isolated on every generic bilinear Hessian
ray inside the cone of nonnegatively curved metrics.
\end{corollary}

\begin{proof}
Having a simple largest singular value is an open dense condition on
matrices.  At $\gamma=0$ the metric is the nonnegatively curved round
product.  For every other Riemannian parameter,
Theorem~\ref{thm:spectral-gap} supplies a plane of strictly negative
sectional curvature.
\end{proof}

The spectral-gap theorem therefore leaves the repeated-top stratum as the
only possible source of a nontrivial nonnegatively curved metric.  At its
maximally symmetric endpoint on $\Sph^2\times\Sph^2$, the mixed
curvature sign reverses completely: no background-mixed plane is
negative.

\begin{theorem}[Isotropic mixed-curvature nonnegativity]
\label{thm:isotropic-mixed}
Let $m=n=2$ and $A=\sigma Q$, where $\sigma>0$ and $Q\in O(3)$.
Put $\delta=\gamma\sigma$, so the Riemannian interval is
$|\delta|<1/2$.  Fix a point $(x,y)$ and background-unit tangent
vectors $X\perp x$ and $Y\perp y$.  Complete them to orthonormal
frames
\[
 \mathcal E=(x,X,\widehat X),\qquad
 \mathcal F=(y,Y,\widehat Y)
\]
whose orientations are chosen so that
$R=\mathcal E^{\mathsf T}Q\mathcal F\in SO(3)$.  Represent $R$ by a
unit quaternion
\[
 q_0+q_1\mathbf i+q_2\mathbf j+q_3\mathbf k
\]
and set
\[
 \omega=q_0^2,\qquad \xi=q_1^2,\qquad
 \upsilon=q_2^2,\qquad \zeta=q_3^2.
\]
Then $\omega+\xi+\upsilon+\zeta=1$ and
\begin{equation}\label{eq:isotropic-curvature}
 \secop_{g_{\gamma,A}}(X\wedge Y)
 =\frac{2\delta^2(1-2\delta R_{00})}
        {\det(G_{\delta,R})\,\Delta_{\delta,R}(X,Y)}
   \mathcal S_\delta,
\end{equation}
where $G_{\delta,R}$ is the matrix of $g_{\gamma,A}$ in the adapted
tangent frame, $\Delta_{\delta,R}(X,Y)$ is the squared area of the
plane, and
\begin{align}
 \mathcal S_\delta={}&
 \xi\zeta(1-2\delta\xi)(1+2\delta\zeta)
 +\omega\upsilon(1-2\delta\omega)(1+2\delta\upsilon)
 \notag\\
 &+8\delta^2\omega\xi\upsilon\zeta.       \label{eq:isotropic-sos}
\end{align}
Consequently every background-mixed sectional curvature is
nonnegative throughout the Riemannian interval.  If $\delta\neq0$,
equality holds precisely when
\[
 \xi\zeta=0,\qquad \omega\upsilon=0.
\]
\end{theorem}

\begin{proof}
Scaling $A$ replaces $\gamma$ by $\delta$, so it is enough to work
with $Q$.  Relative to the adapted frames, write
\[
 R=\begin{pmatrix}
 s&\beta&b_\perp\\
 \alpha&p&q\\
 a_\perp&r&z
 \end{pmatrix},
 \qquad
 C=\begin{pmatrix}p&q\\r&z\end{pmatrix},
 \qquad \lambda=1-\delta s.
\]
The tangent metric has matrix
\begin{equation}\label{eq:isotropic-G}
 G_{\delta,R}=
 \begin{pmatrix}\lambda I_2&\delta C\\
                 \delta C^{\mathsf T}&\lambda I_2
 \end{pmatrix}.
\end{equation}
The three covectors in Theorem~\ref{thm:mixed-curvature} are
\[
 \eta_{XX}=\tfrac12(-\alpha,a_\perp,-\beta,-b_\perp),\quad
 \eta_{XY}=\tfrac12(-\beta,0,-\alpha,0),
\]
\[
 \eta_{YY}=\tfrac12(-\alpha,-a_\perp,-\beta,b_\perp).
\]
A direct block inversion gives
\begin{equation}\label{eq:isotropic-N-over-D}
 G_{\delta,R}^{-1}(\eta_{XY},\eta_{XY})
 -G_{\delta,R}^{-1}(\eta_{XX},\eta_{YY})
 =\frac{N}{4D},
\end{equation}
where $D=\det(G_{\delta,R})>0$ and
\begin{align*}
 N={}&a_\perp^2\lambda
       \bigl(\lambda^2-\delta^2(p^2+q^2)\bigr)
 +b_\perp^2\lambda
       \bigl(\lambda^2-\delta^2(p^2+r^2)\bigr)\\
 &+2a_\perp b_\perp\delta
       \bigl(\lambda^2z-\delta^2p(pz-qr)\bigr)\\
 &+\lambda\delta^2(\beta^2-\alpha^2)(q^2-r^2).
\end{align*}

Use the standard quaternion matrix
\[
 R=\begin{pmatrix}
 1-2(q_2^2+q_3^2)&2(q_1q_2-q_0q_3)&2(q_1q_3+q_0q_2)\\
 2(q_1q_2+q_0q_3)&1-2(q_1^2+q_3^2)&2(q_2q_3-q_0q_1)\\
 2(q_1q_3-q_0q_2)&2(q_2q_3+q_0q_1)&1-2(q_1^2+q_2^2)
 \end{pmatrix}.
\]
Substitution into the preceding numerator, followed only by the unit
relation $q_0^2+q_1^2+q_2^2+q_3^2=1$, gives
\begin{equation}\label{eq:isotropic-factorization}
 N=8(1-2\delta R_{00})\mathcal S_\delta.
\end{equation}
Equations~\eqref{eq:exact-mixed-curvature},
\eqref{eq:isotropic-N-over-D}, and
\eqref{eq:isotropic-factorization} prove
\eqref{eq:isotropic-curvature}.

It remains only to read the sign.  Since $|\delta|<1/2$ and each of
$\omega,\xi,\upsilon,\zeta$ lies in $[0,1]$, every linear factor in
\eqref{eq:isotropic-sos} is positive.  Also
$1-2\delta R_{00}>0$, while $D$ and the squared area are positive by
the Riemannian assumption.  Thus the curvature is nonnegative.  For
$\delta\neq0$, the sum vanishes exactly when its first two products
vanish, which is the stated equality condition.
\end{proof}

The oblique obstruction propagates to arbitrary dimensions by a
fixed-point reduction that will also underlie the flat-torus
construction.

\begin{proposition}[Singular-subsphere reduction]
\label{prop:singular-subsphere}
Choose a singular frame
\[
 Av_i=\sigma_i u_i,
 \qquad A^{\mathsf T}u_i=\sigma_i v_i,
 \qquad 1\leq i\leq q,
\]
and let $I\subset\{1,\ldots,q\}$ have cardinality $r\geq2$.  Set
\[
 N_I=
 \bigl(\Sph^m\cap\spanop\{u_i:i\in I\}\bigr)
 \times
 \bigl(\Sph^n\cap\spanop\{v_i:i\in I\}\bigr).
\]
Whenever $g_{\gamma,A}$ is Riemannian, $N_I\cong
\Sph^{r-1}\times\Sph^{r-1}$ is totally geodesic.  Its induced metric
is precisely the Hessian family associated to the diagonal restriction
\[
 A_I:\spanop\{v_i:i\in I\}\longrightarrow
       \spanop\{u_i:i\in I\},
 \qquad A_Iv_i=\sigma_i u_i.
\]
Consequently every sectional-curvature obstruction in a singular
submodel is also an obstruction in the ambient product.
\end{proposition}

\begin{proof}
Let $R_I$ be $+I$ on $\spanop\{u_i:i\in I\}$ and $-I$ on its
orthogonal complement, and define $S_I$ analogously using the $v_i$.
The singular decomposition gives
\[
 R_I^{\mathsf T}AS_I=A.
\]
Thus $(x,y)\mapsto(R_Ix,S_Iy)$ preserves $h$ and $s_A$, hence also
$g_{\gamma,A}$.  Its fixed-point set is $N_I$, so the fixed-point
lemma for isometries makes $N_I$ totally geodesic.  The product of
great subspheres is also totally geodesic for $h$.  Therefore the
restriction of $\Hess_hs_A$ is the intrinsic Hessian of
$s_A|_{N_I}$, which is the bilinear potential determined by $A_I$.
The final assertion follows from the Gauss equation.
\end{proof}

The isotropic theorem concerns planes mixed with respect to the
background product splitting.  The next calculation gives a negative
oblique plane at the top singular critical point for every matrix of
rank at least two, not only on the repeated-top stratum.

\begin{proposition}[Top-critical oblique-plane obstruction]
\label{prop:top-critical-oblique}
Let $m,n\geq2$, let $\gamma>0$, and abbreviate
\[
 x=\gamma\sigma_1,\qquad y=\gamma\sigma_2,\qquad
 z=\gamma\sigma_3,\qquad d=1-x.
\]
Choose ordered singular pairs $Av_i=\sigma_i u_i$ and set
\[
 \lambda_i^\pm=d\pm\gamma\sigma_i,
 \qquad \Delta_i=\lambda_i^+\lambda_i^-,
 \qquad i=2,3.
\]
Throughout the Riemannian range $x+y<1$, the four vectors at
$(u_1,v_1)$
\[
 E_i^+=\frac{(u_i,v_i)}{\sqrt{2\lambda_i^+}},
 \qquad
 E_i^-=\frac{(u_i,-v_i)}{\sqrt{2\lambda_i^-}},
 \qquad i=2,3,
\]
are $g_{\gamma,A}$-orthonormal.  The plane spanned by
\begin{equation}\label{eq:top-critical-plane}
 U=\frac{E_2^+-E_2^-}{\sqrt2},
 \qquad
 V=\frac{E_3^++E_3^-}{\sqrt2}
\end{equation}
has unit squared area and
\begin{equation}\label{eq:top-critical-curvature}
 \secop_{g_{\gamma,A}}(U\wedge V)
 =-\frac{\mathcal N(x,y,z)}{2\Delta_2\Delta_3},
\end{equation}
where
\begin{equation}\label{eq:top-critical-numerator}
 \mathcal N(x,y,z)
 =d(y^2+z^2)+(2x-1)
   \left(d^2-\sqrt{\Delta_2\Delta_3}\right).
\end{equation}
If $\sigma_2>0$, this curvature is strictly negative.  The same
conclusion holds for $\gamma<0$ after applying the antipodal map on
the first sphere.

On the repeated-top stratum $x=y=\delta$ and $z=\tau\delta$, put
\[
 a=1-2\delta,
 \qquad \Pi=(1-\delta)^2-\tau^2\delta^2.
\]
Then~\eqref{eq:top-critical-curvature} becomes
\begin{equation}\label{eq:repeated-top-curvature}
 \secop_{g_{\gamma,A}}(U\wedge V)
 =\frac{(1-\delta)(\Pi-\delta)-a\sqrt{a\Pi}}
        {2a\Pi}.
\end{equation}
In particular, at the isotropic endpoint $\tau=1$ it is
$-\delta^2/(2(1-2\delta)^2)$.
\end{proposition}

\begin{proof}
Proposition~\ref{prop:singular-subsphere} reduces the computation to
the totally geodesic top-three singular
$\Sph^2\times\Sph^2$.  At $(u_1,v_1)$ the tangent metric splits into
symmetric and antisymmetric modes with coefficients
$\lambda_i^+$ and $\lambda_i^-$, proving orthonormality.
Direct differentiation in orthographic coordinates gives unit area
and~\eqref{eq:top-critical-curvature}--
\eqref{eq:top-critical-numerator}.

It remains to read the sign.  Put
\[
 S=y^2+z^2,
 \qquad L=\sqrt{(d^2-y^2)(d^2-z^2)}.
\]
The Riemannian inequality gives $d>y\geq z\geq0$.  Hence
$0<L\leq d^2$.  If $x\geq1/2$, then
$\mathcal N\geq dS>0$.  If $x<1/2$, the elementary inequality
\[
 1-\sqrt{(1-r)(1-s)}\leq r+s,
 \qquad 0\leq r,s<1,
\]
applied to $r=y^2/d^2$ and $s=z^2/d^2$ gives
$d^2-L\leq S$.  Consequently
\[
 \mathcal N=dS-(1-2x)(d^2-L)\geq xS>0.
\]
Since $S>0$ exactly when $y>0$, the strict sign follows.  The map
$(x_0,y_0)\mapsto(-x_0,y_0)$ pulls $g_{\gamma,A}$ back to
$g_{-\gamma,A}$, proving the negative-parameter assertion.
Substitution of $x=y=\delta$ and $z=\tau\delta$ gives the final two
formulas.
\end{proof}

\begin{proposition}[Rank-one zero-critical obstruction]
\label{prop:rank-one-oblique}
Assume $m,n\geq2$ and $A$ has rank one.  Write its only positive
singular value as $\sigma$, choose singular pairs so that
$\sigma_1=\sigma$ and $\sigma_2=\sigma_3=0$, and put
$\delta=\gamma\sigma$.  At the zero critical point
$p=(u_2,v_2)$ define
\[
 F_1^\pm=\frac{(u_1,\pm v_1)}{\sqrt{2(1\pm\delta)}},
 \qquad
 F_3^\pm=\frac{(u_3,\pm v_3)}{\sqrt2}.
\]
For $|\delta|<1$, the vectors
\[
 U=\frac{F_1^+-F_1^-}{\sqrt2},
 \qquad
 V=\frac{F_3^++F_3^-}{\sqrt2}=(u_3,0)
\]
are $g_{\gamma,A}$-orthonormal and
\begin{equation}\label{eq:rank-one-oblique}
 \secop_{g_{\gamma,A}}(U\wedge V)
 =-\frac{1-\sqrt{1-\delta^2}}
          {2\sqrt{1-\delta^2}}.
\end{equation}
Thus this plane is strictly negative for every nonzero Riemannian
parameter, and
\begin{equation}\label{eq:rank-one-quadratic}
 \secop_{g_{\gamma,A}}(U\wedge V)
 =-\frac{\gamma^2\sigma^2}{4}+O(\gamma^4).
\end{equation}
\end{proposition}

\begin{proof}
Repeat the orthographic-coordinate calculation in
Proposition~\ref{prop:top-critical-oblique} using the local singular
coordinates $(0,\sigma,0)$ at the critical pair $(u_2,v_2)$.
In the notation of~\eqref{eq:top-critical-numerator}, this means
$x=z=0$, $y=\delta$, and $d=1$.  The squared area is one and the
curvature is
\[
 -\frac{\delta^2-1+\sqrt{1-\delta^2}}
        {2(1-\delta^2)}
 =-\frac{1-\sqrt{1-\delta^2}}
        {2\sqrt{1-\delta^2}},
\]
which is negative for $0<|\delta|<1$.  Its Taylor expansion gives
\eqref{eq:rank-one-quadratic}.
\end{proof}

\begin{corollary}[Four-dimensional detection and complete isolation]
\label{cor:complete-isolation}
Let $m,n\geq2$.  For every nonzero
$A:\R^{n+1}\to\R^{m+1}$, choose an ordered singular frame and set
$N_{123}=N_{\{1,2,3\}}$.  For every nonzero Riemannian parameter,
$N_{123}\cong\Sph^2\times\Sph^2$ contains a plane of strictly
negative ambient sectional curvature.  Consequently,
\[
 \left\{\gamma:
  g_{\gamma,A}\text{ is Riemannian and }
  \secop_{g_{\gamma,A}}\geq0
 \right\}=\{0\}.
\]
\end{corollary}

\begin{proof}
If $\sigma_2>0$, Proposition~\ref{prop:top-critical-oblique} supplies
a negative oblique plane in $N_{123}$.  If $\sigma_2=0$, then $A$ has
rank one and Proposition~\ref{prop:rank-one-oblique} supplies a
negative oblique plane in the same submanifold.  At $\gamma=0$ the
metric is the nonnegatively curved round product.
\end{proof}

\begin{corollary}[Explicit dimension-free negative margin]
\label{cor:uniform-negative-margin}
Fix $0<\epsilon<1/2$, and let $r_\epsilon\in(0,1)$ be the unique
solution of
\begin{equation}\label{eq:epsilon-root}
 (1-r_\epsilon)^2=\epsilon r_\epsilon^3.
\end{equation}
Set
\[
 c_\epsilon=\frac{\epsilon^4r_\epsilon^3}{16}
 >\frac{\epsilon^4}{128}.
\]
For every $m,n\geq2$, every nonzero $A$, and every $\gamma$ satisfying
\[
 \epsilon\leq |\gamma|(\sigma_1(A)+\sigma_2(A))
 \leq1-\epsilon,
\]
one has
\[
 \min_{\Pi\in\operatorname{Gr}_2(T(\Sph^m\times\Sph^n))}
 \secop_{g_{\gamma,A}}(\Pi)\leq-c_\epsilon.
\]
\end{corollary}

\begin{proof}
The antipodal isometry lets us assume $\gamma>0$.  Put
\[
 x=\gamma\sigma_1,\qquad y=\gamma\sigma_2,
 \qquad z=\gamma\sigma_3,\qquad p=x+y,\qquad q=\frac yp.
\]
Thus $0\leq q\leq1/2$ and $\epsilon\leq p\leq1-\epsilon$.
When $x>y$, Theorem~\ref{thm:spectral-gap}, together with
$z\leq y$ and the fact that its last denominator is at most one,
gives a plane $\Pi_{\mathrm{gap}}$ with
\begin{align}
 -\secop(\Pi_{\mathrm{gap}})
 &\geq\frac{(x^2-y^2)^3}{16(x^2+y^2)}\notag\\
 &\geq\frac{\epsilon^4(1-2q)^3}{16}.
 \label{eq:quantitative-gap}
\end{align}
When $x=y$, define the first lower bound to be zero; only the oblique
estimate below is then used.

For the oblique plane in Proposition~\ref{prop:top-critical-oblique},
its proof gives
\[
 \mathcal N(x,y,z)\geq\min\{x,1-x\}(y^2+z^2).
\]
Moreover $\Delta_2\Delta_3\leq(1-x)^4\leq1$,
$x\geq p/2\geq\epsilon/2$, and
$1-x\geq1-p\geq\epsilon$.  Therefore
\begin{equation}\label{eq:quantitative-oblique}
 -\secop(\Pi_{\mathrm{obl}})
 \geq\frac{\epsilon^3q^2}{4}.
\end{equation}
Writing $r=1-2q$, the larger of
\eqref{eq:quantitative-gap} and~\eqref{eq:quantitative-oblique} is at
least
\[
 \frac1{16}\max\{\epsilon^4r^3,
                         \epsilon^3(1-r)^2\}.
\]
The increasing and decreasing terms meet exactly at the unique root
in~\eqref{eq:epsilon-root}; hence their maximum is at least
$c_\epsilon$.  Finally $r_\epsilon>1/2$, because at $r=1/2$ the two
sides of~\eqref{eq:epsilon-root} are respectively $1/4$ and
$\epsilon/8$.  This proves the simpler displayed bound as well.
\end{proof}

The following specialization uses a different plane than the canonical
choice in Theorem~\ref{thm:spectral-gap} and gives a particularly simple
curvature function.

\begin{corollary}[A simple zero-threshold specialization]
\label{cor:zero-threshold}
On $\Sph^2\times\Sph^2$, let
\[
 A_0=\begin{pmatrix}0&2&0\\1&0&0\\0&1&0\end{pmatrix}.
\]
At $p=(e_0,e_0)$ take the background-mixed plane spanned by
\[
 \mathsf X=((0,1,0),0),\qquad
 \mathsf Y=(0,(0,1,0)).
\]
For every
\[
 0<\lvert\gamma\rvert<\frac1{\sqrt5+1},
\]
the metric $g_{\gamma,A_0}$ is Riemannian on the whole product and
\begin{equation}\label{eq:negative-all-small}
 \secop_{g_{\gamma,A_0}}(\mathsf X\wedge\mathsf Y)
 =-\frac{3\gamma^4}{4(1-\gamma^2)}<0.
\end{equation}
Thus no matrix-independent positive interval of mixed-curvature
nonnegativity exists, even after normalizing by the sharp Riemannian
endpoint.
\end{corollary}

\begin{proof}
Since $A_0^{\mathsf T}A_0=\operatorname{diag}(1,5,0)$, the Riemannian
interval follows
from Proposition~\ref{prop:positive}.  At $p$ one has
\[
 a=(1,0),\qquad b=(2,0),\qquad
 P_xA_0P_y=\begin{pmatrix}0&0\\1&0\end{pmatrix}.
\]
In the ordered tangent basis the metric and the three covectors from
Theorem~\ref{thm:mixed-curvature} are
\[
 g_{\gamma,A_0}=
 \begin{pmatrix}
  1&0&0&0\\0&1&\gamma&0\\0&\gamma&1&0\\0&0&0&1
 \end{pmatrix}.
\]
\[
 \eta_{XX}=\eta_{YY}=(-1/2,0,-1,0),
 \quad
 \eta_{XY}=(-1,0,-1/2,0).
\]
Formula~\eqref{eq:exact-mixed-curvature} now reduces to
\eqref{eq:negative-all-small}; the squared area is one.
\end{proof}

A different ray illustrates that an individual mixed plane can have a
nonzero positive-to-negative crossing, even though
Theorem~\ref{thm:spectral-gap} already supplies another plane with
negative curvature along the punctured ray.

\begin{proposition}[A plane-wise nonzero sign-change threshold]
\label{prop:positive-threshold}
Let
\[
 B=\frac1{100}
 \begin{pmatrix}1&2&-32\\-48&-23&3\\-39&-15&-3\end{pmatrix},
 \qquad
 \gamma_{\mathrm R}=\frac1{\sigma_1(B)+\sigma_2(B)}.
\]
At $p=(e_0,e_0)$ consider the plane spanned by
\[
 \mathsf X=\big((0,63/100,78/100),0\big),\qquad
 \mathsf Y=\big(0,(0,9/100,-1)\big).
\]
There is a unique $\gamma_*\in(0,\gamma_{\mathrm R})$ at which its
sectional curvature changes sign, and
\begin{equation}\label{eq:threshold-bracket}
 \frac{4679671}{5000000}<\gamma_*
 <\frac{9359343}{10000000},
 \qquad
 \gamma_*=0.935934207005572\ldots.
\end{equation}
The curvature is positive for $0<\gamma<\gamma_*$, zero at
$\gamma_*$, and negative for $\gamma_*<\gamma<\gamma_{\mathrm R}$.
\end{proposition}

\begin{proof}
Exact substitution in~\eqref{eq:exact-mixed-curvature} gives
\begin{equation}\label{eq:threshold-rational-function}
 \secop_{g_{\gamma,B}}(\mathsf X\wedge\mathsf Y)
 =\frac{11\gamma^2P(\gamma)}{4Q(\gamma)L(\gamma)},
\end{equation}
where
\begin{align*}
 P(\gamma)={}&1873383134707\gamma^3
 -187783301783700\gamma^2\\
 &-6938955210000\gamma+169451579000000,\\
 Q(\gamma)={}&28151192500-563023850\gamma-7287743\gamma^2,\\
 L(\gamma)={}&489\gamma^4+6160\gamma^3-306400\gamma^2
                    -160000\gamma+4000000.
\end{align*}
On the Riemannian interval, $Q$ and $L$ are positive: in fact, up to
positive constants they are respectively the squared area of this
plane and the determinant of $g_{\gamma,B}$ at $p$.

The crossing lies strictly inside the global Riemannian interval.
Indeed,
\[
 \left(\frac{47}{50}\right)^2 2\norm{B}_{\mathrm F}^2
 =\frac{6213917}{6250000}<1,
\]
so $\gamma_{\mathrm R}>47/50$.  On the other hand,
$\sigma_1(B)+\sigma_2(B)\geq\norm{Be_0}
=\sqrt{3826}/100>1/2$, so $\gamma_{\mathrm R}<2$.

On $[0,2]$ the cubic $P$ is strictly decreasing.  Indeed,
\[
 \begin{aligned}
  P''(\gamma)&=11240298808242\gamma-375566603567400<0,\\
  P'(0)&=-6938955210000<0.
 \end{aligned}
\]
Finally, exact rational substitution gives
\[
 P(4679671/5000000)>0,\qquad
 P(9359343/10000000)<0.
\]
Thus $P$ has exactly one zero in the Riemannian interval, with the
bracket~\eqref{eq:threshold-bracket}; the sign statement follows from
\eqref{eq:threshold-rational-function}.
\end{proof}

\begin{remark}
Corollary~\ref{cor:zero-threshold} shows that the nonnegative
$\gamma^2$ coefficient in~\eqref{eq:mixed-expansion} is not a uniform
curvature barrier.  When $X$ is parallel to $a$ and $Y$ is parallel to
$b$, that coefficient vanishes and the first nonzero term can be
negative; in~\eqref{eq:negative-all-small} it is
$-3\gamma^4/4+O(\gamma^6)$.
\end{remark}

\section{Flat totally geodesic tori}\label{sec:tori}

Fix a singular value decomposition
\begin{equation}\label{eq:svd}
 Av_k=\sigma_k u_k,
 \qquad 1\leq k\leq q,
\end{equation}
with orthonormal bases completed in the left and right null spaces if
necessary.  For $i<j$, let
\begin{equation}\label{eq:tori}
 C_{ij}=\Sph^m\cap\spanop\{u_i,u_j\},\qquad
 D_{ij}=\Sph^n\cap\spanop\{v_i,v_j\},\qquad
 T_{ij}=C_{ij}\times D_{ij}.
\end{equation}
Each $T_{ij}$ is a smoothly embedded $2$-torus.

\begin{proposition}\label{prop:tg}
Whenever $g_{\gamma,A}$ is Riemannian, every torus $T_{ij}$ in
\eqref{eq:tori} is totally geodesic for $g_{\gamma,A}$.
\end{proposition}

\begin{proof}
This is Proposition~\ref{prop:singular-subsphere} with
$I=\{i,j\}$.  The argument also covers the case where one of the
paired great circles is the entire sphere.
\end{proof}

\begin{proposition}\label{prop:flat}
The metric induced by $g_{\gamma,A}$ on every $T_{ij}$ is flat.
\end{proposition}

\begin{proof}
Parametrize the two great circles by
\[
 x(u)=\cos u\,u_i+\sin u\,u_j,
 \qquad
 y(v)=\cos v\,v_i+\sin v\,v_j.
\]
The restriction of $h$ is $du^2+dv^2$.  Since $T_{ij}$ is also
totally geodesic for $h$, the ambient Hessian restricts to the
intrinsic Hessian of
\begin{equation}\label{eq:phi-uv}
 \phi(u,v)=\restr{s_A}{T_{ij}}
 =\sigma_i\cos u\cos v+\sigma_j\sin u\sin v.
\end{equation}
Introduce local coordinates
\[
 \alpha=u+v,\qquad \beta=u-v,
\]
and abbreviate
\begin{equation}\label{eq:ab}
 a=\frac{\sigma_i+\sigma_j}{2},
 \qquad b=\frac{\sigma_i-\sigma_j}{2}.
\end{equation}
The elementary product-to-sum identities give
\[
 \phi=b\cos\alpha+a\cos\beta,
 \qquad
 du^2+dv^2=\frac12(d\alpha^2+d\beta^2).
\]
It follows that
\begin{equation}\label{eq:separated-metric}
 \restr{g_{\gamma,A}}{T_{ij}}
 =\left(\frac12-\gamma b\cos\alpha\right)d\alpha^2
  +\left(\frac12-\gamma a\cos\beta\right)d\beta^2.
\end{equation}
Both coefficients are positive because the ambient metric is positive
definite.  Each depends on only its own coordinate.  Locally making
the changes of variables
\[
 \widehat\alpha=\int_0^\alpha
 \sqrt{\frac12-\gamma b\cos t}\,dt,
 \qquad
 \widehat\beta=\int_0^\beta
 \sqrt{\frac12-\gamma a\cos t}\,dt
\]
turns~\eqref{eq:separated-metric} into
$d\widehat\alpha^2+d\widehat\beta^2$.  Hence its Gaussian curvature
vanishes identically.
\end{proof}

\begin{proof}[Proof of Theorem~\ref{thm:main}]
Proposition~\ref{prop:positive} gives the exact Riemannian range.  By
Propositions~\ref{prop:tg} and~\ref{prop:flat}, every chosen singular
frame supplies $\binom q2$ flat totally geodesic tori.  The Gauss
equation for a totally geodesic surface identifies
its intrinsic Gaussian curvature with the ambient sectional curvature
of its tangent plane.  That curvature is therefore zero at every
point of every $T_{ij}$, which rules out strictly positive sectional
curvature.  When $m,n\geq2$ and $\sigma_2>0$,
Proposition~\ref{prop:top-critical-oblique} supplies strict negative
curvature on the totally geodesic singular
$\Sph^2\times\Sph^2$.  In the remaining rank-one case,
Proposition~\ref{prop:rank-one-oblique} supplies strict negative
curvature instead.  The assertion for $A=0$ is immediate from the
definition.
\end{proof}

\begin{remark}[Repeated and zero singular values]
The theorem does not require $A$ to have full rank or simple singular
values.  When a positive singular value is repeated, the choice of
singular frame is nonunique and may produce continuous families of
flat tori.  When some $\sigma_j=0$, the same construction still works:
the corresponding paired great circles are fixed by the involution
and the separated formula~\eqref{eq:separated-metric} remains valid.
The count $\binom q2$ is therefore a guaranteed count for each chosen
frame, not a classification of all flat tori.
\end{remark}

\section{Inherited symmetry and the full isometry group}\label{sec:symmetry}

The flat-torus proof uses finite involutions, but it is useful to state
exactly how much of the background product symmetry is preserved.  Let
\begin{equation}\label{eq:GA}
 G_A=\big\{(R,S)\in O(m+1)\times O(n+1):
 R^{\mathsf T}AS=A\big\}.
\end{equation}
This is the stabilizer of $A$ for the natural left-right action.

\begin{proposition}\label{prop:symmetry}
For $\gamma\neq0$, the factor-preserving isometries of the round
product which remain isometries of $g_{\gamma,A}$ are exactly $G_A$.
If the distinct positive singular values of $A$ have multiplicities
$d_1,\ldots,d_k$ and $r_A=\rank A$, then
\begin{equation}\label{eq:GA-structure}
 G_A\cong
 \left(\prod_{\ell=1}^k O(d_\ell)_{\mathrm{diag}}\right)
 \times O(m+1-r_A)\times O(n+1-r_A).
\end{equation}
Here $O(d_\ell)_{\mathrm{diag}}$ acts by the same orthogonal matrix on
the corresponding left and right singular subspaces.
\end{proposition}

\begin{proof}
If $(R,S)\in G_A$, then
\[
 s_A(Rx,Sy)=\ip{x}{R^{\mathsf T}ASy}=s_A(x,y).
\]
The associated product isometry therefore preserves $h$ and
$\Hess_hs_A$, so it preserves $g_{\gamma,A}$.

Conversely, suppose the product isometry associated to $(R,S)$
preserves $g_{\gamma,A}$.  Since it already preserves $h$ and
$\gamma\neq0$, the function
\[
 f=s_{R^{\mathsf T}AS}-s_A
\]
satisfies $\Hess_h f=0$.  On a compact connected manifold such a
function is constant: at a maximum its gradient vanishes, while
$\Hess f=0$ makes the gradient parallel.  Both bilinear functions have
zero integral over the product of spheres, so the constant is zero.
Thus $R^{\mathsf T}AS=A$.

Finally, write $A$ in singular-value block form.  The relation
$R^{\mathsf T}AS=A$ forces $R$ and $S$ to preserve the eigenspaces of
$AA^{\mathsf T}$ and $A^{\mathsf T}A$.  On a positive singular block
of value $\rho$, the relation is
$R_\rho^{\mathsf T}(\rho I)S_\rho=\rho I$, hence
$R_\rho=S_\rho$.  On the two null spaces the orthogonal actions are
independent.  This gives~\eqref{eq:GA-structure}.
\end{proof}

\begin{corollary}[Curvature mechanisms and symmetry]
\label{cor:curvature-symmetry}
Assume $m,n\geq2$, $A\neq0$, and $\gamma\neq0$ is in the Riemannian
interval.  Then one of the following holds:
\begin{enumerate}
 \item $\sigma_1>\sigma_2$, and $g_{\gamma,A}$ has a background-mixed
 plane of strictly negative sectional curvature;
 \item $\sigma_1=\sigma_2$, and $g_{\gamma,A}$ has a negative oblique
 plane in a totally geodesic singular $\Sph^2\times\Sph^2$, while the
 inherited factor-preserving isometry group $G_A$ contains a circle
 subgroup.
\end{enumerate}
\end{corollary}

\begin{proof}
The first alternative is Theorem~\ref{thm:spectral-gap}.  In the
second, Proposition~\ref{prop:top-critical-oblique} gives the negative
plane.  The top positive singular value has multiplicity at least two,
so Proposition~\ref{prop:symmetry} also exhibits a diagonal $O(d)$
factor with $d\geq2$.  Its identity component contains a circle.
\end{proof}

\begin{corollary}[The generic $\Sph^2\times\Sph^2$ case]
If $A:\R^3\to\R^3$ is invertible with three distinct singular values,
then the inherited factor-preserving group is
\[
 G_A\cong O(1)^3\cong(\mathbb Z/2\mathbb Z)^3.
\]
It contains the three involutions whose fixed sets are the tori
$T_{12}$, $T_{13}$, and $T_{23}$, but it contains no circle subgroup.
\end{corollary}

\begin{remark}[What is not being claimed]
Proposition~\ref{prop:symmetry} computes the intersection of
$\Isom(M,g_{\gamma,A})$ with the factor-preserving isometry group of
the round product.  It does not classify the full isometry group of
$g_{\gamma,A}$.  For example, equal-dimensional factors can admit a
factor interchange under an additional algebraic condition, and an
exceptional parameter could in principle have isometries not inherited
from $h$.  Accordingly, the finite group in the corollary is not
asserted to be the full isometry group.  The curvature obstruction does
not require such an assertion.
\end{remark}

\section{A two-jet repair criterion}\label{sec:two-jet}

The exact rays above have a vanishing first-order curvature signal for
a structural reason:
\begin{equation}\label{eq:hessian-lie}
 \Hess_h f=\frac12\mathcal L_{\nabla f}h.
\end{equation}
Thus their first metric variation is tangent to the diffeomorphism
orbit of $h$.  The next statement separates that gauge term from the
first curvature-relevant correction.  It applies to an arbitrary
Riemannian product; we then specialize it to $s_A$.

Let $h=h_1\oplus h_2$ on $M_1\times M_2$, let $\nabla$ be its
connection, and let $B$ be a symmetric $2$-tensor.  For an $h$-unit
mixed pair $X\in TM_1$, $Y\in TM_2$, define
\begin{align}
 \mathfrak m_B(X,Y)=\frac12\big(&
  (\nabla^2_{X,Y}B)(X,Y)+(\nabla^2_{Y,X}B)(X,Y)\notag\\
 &-(\nabla^2_{X,X}B)(Y,Y)
  -(\nabla^2_{Y,Y}B)(X,X)\big).
 \label{eq:mixed-jet-functional}
\end{align}
Because $R^h(X,Y)=0$, the two cross derivatives agree, but the
symmetrized form makes the curvature linearization transparent.

\begin{theorem}[Two-jet normal form and mixed repair condition]
\label{thm:two-jet}
Let $g_t$ be a smooth family of metrics on $M_1\times M_2$ with
\begin{equation}\label{eq:two-jet-path}
 g_t=h+t\Hess_h f+t^2B+O(t^3)
\end{equation}
in $C^2$.  Let $\Phi_t$ be the local flow of
$-\frac12\nabla f$.  Then
\begin{equation}\label{eq:two-jet-normal-form}
 \Phi_t^*g_t=h+t^2\widehat B+O(t^3),
 \qquad
 \widehat B=B-\frac18\mathcal L_{\nabla f}^2h.
\end{equation}
For every $p\in M_1\times M_2$ and every $h$-unit mixed pair
$X,Y$ at $p$,
\begin{equation}\label{eq:two-jet-curvature}
 \secop_{g_t}\bigl(d\Phi_t(X\wedge Y)\bigr)
 =t^2\mathfrak m_{\widehat B}(X,Y)+O(t^3).
\end{equation}
Consequently, a necessary condition for $g_t$ to have nonnegative
sectional curvature for all sufficiently small $t$ is
\begin{equation}\label{eq:two-jet-necessary}
 \mathfrak m_{\widehat B}(X,Y)\geq0
 \quad\text{for every background-mixed unit pair }X,Y.
\end{equation}
If the inequality is strict in the opposite direction for one pair,
$g_t$ has negative sectional curvature for every sufficiently small
nonzero $t$.
\end{theorem}

\begin{proof}
Put $Z=\nabla f$.  Pullback by the flow of $-Z/2$ acts to second order
as $\exp(-\frac t2\mathcal L_Z)$.  Using
\eqref{eq:hessian-lie} in~\eqref{eq:two-jet-path}, the coefficient of
$t$ cancels, while the coefficient of $t^2$ is
\[
 B-\frac12\mathcal L_Z\Hess f
   +\frac18\mathcal L_Z^2h
 =B-\frac18\mathcal L_Z^2h.
\]
This proves~\eqref{eq:two-jet-normal-form}.

Choose product $h$-normal coordinates at $p$ whose first coordinate
directions include $X$ and $Y$.  The lower-index curvature formula for
$h+t^2\widehat B+O(t^3)$, evaluated on $(X,Y,X,Y)$, has coefficient
\[
 \frac12\bigl(
  (\nabla^2_{X,Y}\widehat B)(X,Y)
 +(\nabla^2_{Y,X}\widehat B)(X,Y)
 -(\nabla^2_{X,X}\widehat B)(Y,Y)
 -(\nabla^2_{Y,Y}\widehat B)(X,X)
 \bigr).
\]
All background curvature terms with both factor directions vanish,
as do the Christoffel-quadratic terms at this order.  The squared area
of the plane is $1+O(t^2)$, so division gives
\eqref{eq:two-jet-curvature}.  Sectional curvature is preserved by
pullback, and the sign conclusions follow.
\end{proof}

On a closed product the inequalities in
\eqref{eq:two-jet-necessary} satisfy a global compatibility identity.
The functional~\eqref{eq:mixed-jet-functional} is the classical
mixed-plane curvature linearization
\cite{BourguignonDeschampsSentenac1973}; tracing it gives the following
global consequence in the present gauge-normalized setting.
Let $e_1,\ldots,e_m$ and $f_1,\ldots,f_n$ be local product
orthonormal frames, and write $\tr_1B$ and $\tr_2B$ for the partial
traces of $B$.

\begin{theorem}[Integrated mixed two-jet rigidity]
\label{thm:integrated-two-jet}
If $M_1$ and $M_2$ are closed, then every symmetric $2$-tensor $B$
satisfies
\begin{equation}\label{eq:integrated-two-jet}
 \int_{M_1\times M_2}
 \sum_{i=1}^m\sum_{\alpha=1}^n
 \mathfrak m_B(e_i,f_\alpha)\,dV_h=0.
\end{equation}
Consequently, if
$\mathfrak m_B(X,Y)\geq0$ for every background-mixed unit pair, then
\begin{equation}\label{eq:mixed-two-jet-zero}
 \mathfrak m_B(X,Y)=0
 \quad\text{for every background-mixed pair }X,Y.
\end{equation}
In particular, any family in Theorem~\ref{thm:two-jet} which has
nonnegative sectional curvature for all sufficiently small $t$ must
satisfy
\begin{equation}\label{eq:normalized-two-jet-zero}
 \mathfrak m_{\widehat B}(X,Y)=0
 \quad\text{for every background-mixed pair }X,Y.
\end{equation}
Thus every such mixed curvature is $O(t^3)$ after gauge
normalization; none can be lifted strictly at quadratic order without
another becoming negative.
\end{theorem}

\begin{proof}
Tracing~\eqref{eq:mixed-jet-functional} in product normal frames gives
\begin{align}
 \sum_{i,\alpha}\mathfrak m_B(e_i,f_\alpha)
 ={}&\operatorname{div}_1\operatorname{div}_2(B|_{TM_1\otimes TM_2})
 \notag\\
 &-\frac12\Delta_1(\tr_2B)
  -\frac12\Delta_2(\tr_1B).
 \label{eq:two-jet-divergence}
\end{align}
The two possible orders in the first term agree because the mixed
curvature of the product connection vanishes.  Each term on the
right-hand side is a divergence.  Its integral is zero on the closed
product, proving~\eqref{eq:integrated-two-jet}.

Now suppose every mixed value is nonnegative.  Its trace in any
product orthonormal frame is a nonnegative function with zero
integral, hence vanishes identically.  Every summand in that frame is
therefore zero.  At a given point any prescribed mixed unit pair can
be included in such frames, which proves~\eqref{eq:mixed-two-jet-zero}.
The last assertion follows from
Theorem~\ref{thm:two-jet}.
\end{proof}

The global test is therefore feasible, but only at equality.  Put
$Z=\nabla f$, let $\Psi_t$ be the flow of $Z/2$, and set
\begin{equation}\label{eq:gauge-repair}
 B_{\mathrm{gauge}}=\frac18\mathcal L_Z^2h.
\end{equation}
Then
\[
 \Psi_t^*h
 =h+t\Hess f+t^2B_{\mathrm{gauge}}+O(t^3),
 \qquad \widehat B=0.
\]
This correction solves all the two-jet equalities, but produces only
a family isometric to the original product.  The kernel is not purely
gauge: tensors pulled back from the individual factors also satisfy
\eqref{eq:mixed-two-jet-zero}, as expected from ordinary product
variations.  Theorem~\ref{thm:full-coupled-kernel} below classifies the
full kernel modulo these two trivial sectors.  Testing its first
nonzero higher-order curvature is a separate problem.
The higher-order positive mixed variations constructed in
\cite{BourguignonDeschampsSentenac1973} show why this second step is
not vacuous; they also illustrate that improving all mixed planes does
not by itself produce positive sectional curvature on every plane.

Condition~\eqref{eq:normalized-two-jet-zero} is necessary, not
sufficient: when its coefficient vanishes, higher jets can decide the
sign.  It is invariant under the remaining second-order gauge freedom.
Indeed, replacing $\widehat B$ by
$\widehat B+\mathcal L_Wh$ changes the linearized curvature by a Lie
derivative of the product curvature tensor, which vanishes on a mixed
four-tuple.

For the bilinear potentials, the abstract condition gives explicit
values which any proposed nonlinear repair must attain.

\begin{corollary}[Explicit two-jet cost for bilinear Hessian paths]
\label{cor:two-jet-cost}
Assume $m,n\geq2$, let $A\neq0$, and consider
\begin{equation}\label{eq:corrected-bilinear-path}
 g_t=h+t\Hess_hs_A+t^2B+O(t^3).
\end{equation}
Suppose $g_t$ has nonnegative sectional curvature for all sufficiently
small $t$.  Choose an ordered singular frame and pad the singular
values by zeros.

If $\sigma_2>0$, set
\[
 p=(u_1,v_1),\qquad X=(u_3,0),\qquad Y=(0,v_2).
\]
Then nonnegative curvature forces the exact repair identity
\begin{equation}\label{eq:two-jet-cost-rank-two}
 \mathfrak m_B(X,Y)=
 \frac{\sigma_2^2+\sigma_3^2}{4}.
\end{equation}
If $A$ has rank one, set instead
\[
 p=(u_2,v_2),\qquad X=(u_3,0),\qquad Y=(0,v_1).
\]
Then nonnegative curvature forces
\begin{equation}\label{eq:two-jet-cost-rank-one}
 \mathfrak m_B(X,Y)=\frac{\sigma_1^2}{4}.
\end{equation}
In particular, $B=0$ violates one of these identities for every
$A\neq0$.  Hence every nontrivial exact bilinear Hessian ray is
already obstructed by its two-jet modulo diffeomorphisms.
\end{corollary}

\begin{proof}
First suppose $\sigma_2>0$.  At $p=(u_1,v_1)$ the gradient of $s_A$
vanishes, and the Hessian endomorphism has eigenvectors
\[
 e_i^\pm=\frac{(u_i,\pm v_i)}{\sqrt2},
 \qquad
 \Hess s_A(e_i^\pm)=(-\sigma_1\pm\sigma_i)e_i^\pm.
\]
Consequently the differential of the flow in
Theorem~\ref{thm:two-jet}, applied to $Y=(0,v_2)$ and $X=(u_3,0)$,
agrees to first order with the metric-normalized vectors
in~\eqref{eq:top-critical-plane}.  More explicitly,
\[
 d\Phi_t(e_i^\pm)
 =\left(1+\frac t2(\sigma_1\mp\sigma_i)\right)e_i^\pm+O(t^2)
 =(1-t\sigma_1\pm t\sigma_i)^{-1/2}e_i^\pm+O(t^2).
\]
Expanding
\eqref{eq:top-critical-curvature} with
$(x,y,z)=t(\sigma_1,\sigma_2,\sigma_3)$ gives
\begin{equation}\label{eq:top-critical-two-jet}
 \secop=-\frac{t^2}{4}(\sigma_2^2+\sigma_3^2)+O(t^3).
\end{equation}
Theorem~\ref{thm:two-jet} therefore gives
\[
 \mathfrak m_{-\frac18\mathcal L_{\nabla s_A}^2h}(X,Y)
 =-\frac{\sigma_2^2+\sigma_3^2}{4}.
\]
Hence
\[
 \mathfrak m_{\widehat B}(X,Y)
 =\mathfrak m_B(X,Y)-\frac{\sigma_2^2+\sigma_3^2}{4}.
\]
Theorem~\ref{thm:integrated-two-jet} forces this value to vanish,
which proves~\eqref{eq:two-jet-cost-rank-two}.

If $A$ has rank one, the same comparison at $(u_2,v_2)$ uses the
vectors in Proposition~\ref{prop:rank-one-oblique}.  Equation
\eqref{eq:rank-one-quadratic} gives
\[
 \mathfrak m_{-\frac18\mathcal L_{\nabla s_A}^2h}(X,Y)
 =-\frac{\sigma_1^2}{4},
\]
so $\mathfrak m_{\widehat B}(X,Y)
=\mathfrak m_B(X,Y)-\sigma_1^2/4$.  Again
Theorem~\ref{thm:integrated-two-jet} forces this value to vanish,
proving~\eqref{eq:two-jet-cost-rank-one}.  The right-hand side is
positive in both cases.
\end{proof}

\subsection{The full coupled kernel modulo gauge and product tensors}

The partial-divergence reduction below goes back to
Bourguignon--Deschamps--Sentenac~\cite{BourguignonDeschampsSentenac1972}.
Their analysis settles the case without Killing fields and, in the general
case, records the coupled subspace generated by Killing one-forms.  The
tensor-product structure of the remaining equation identifies that subspace
as the entire kernel.

Assume throughout this subsection that $M_1$ and $M_2$ are closed and
connected.  Let
\[
 \mathcal D_i:\Omega^1(M_i)\longrightarrow
 \Gamma(S^2T^*M_i),
 \qquad
 \mathcal D_i\alpha=\frac12\mathcal L_{\alpha^\sharp}h_i
\]
be the Killing operator, and put
\[
 \mathcal K_i=\ker\mathcal D_i,
 \qquad
 \mathcal Z_i=\ker d_i^*\subset\Omega^1(M_i).
\]
Thus $\mathcal K_i$ is the finite-dimensional space of Killing
one-forms, $\mathcal Z_i$ is the space of co-closed one-forms, and
$\mathcal K_i\subset\mathcal Z_i$.  All orthogonal complements below
are taken in $L^2$.  For a mixed covariant tensor
$C\in\Gamma(T^*M_1\boxtimes T^*M_2)$, let $B_C$ denote the symmetric
tensor whose pure blocks vanish and whose mixed block is
\begin{equation}\label{eq:general-mixed-tensor}
 B_C((X,Y),(X',Y'))=C(X,Y')+C(X',Y).
\end{equation}
Finally, define
\begin{align*}
 \mathfrak K_h&=\{B:\mathfrak m_B(X,Y)=0
                     \text{ for every mixed pair }X,Y\},\\
 \mathfrak G_h&=\{\mathcal L_Wh:W\in\Gamma(TM)\},\\
 \mathfrak P_h&=\{\pi_1^*P+\pi_2^*Q:
       P\in\Gamma(S^2T^*M_1),\ Q\in\Gamma(S^2T^*M_2)\}.
\end{align*}

\begin{theorem}[Full coupled-kernel classification]
\label{thm:full-coupled-kernel}
Every element of $\mathfrak K_h$ has a decomposition
\begin{equation}\label{eq:full-kernel-decomposition}
 B=\mathcal L_Wh+\pi_1^*P+\pi_2^*Q+B_C,
\end{equation}
where
\begin{equation}\label{eq:canonical-coupled-space}
 C\in\mathfrak H_{12}:=
 (\mathcal K_1\otimes\mathcal Z_2)
 +(\mathcal Z_1\otimes\mathcal K_2).
\end{equation}
The tensor $B_C$ is uniquely determined by the class of $B$ modulo
$\mathfrak G_h+\mathfrak P_h$.  Equivalently,
\begin{equation}\label{eq:full-kernel-quotient}
 \frac{\mathfrak K_h}{\mathfrak G_h+\mathfrak P_h}
 \cong\mathfrak H_{12}.
\end{equation}
More precisely, the decomposition
\begin{equation}\label{eq:coupled-orthogonal-splitting}
 \mathfrak H_{12}
 = (\mathcal K_1\otimes\mathcal K_2)
 \oplus
   (\mathcal K_1\otimes(\mathcal Z_2\ominus\mathcal K_2))
 \oplus
   ((\mathcal Z_1\ominus\mathcal K_1)\otimes\mathcal K_2)
\end{equation}
is $L^2$-orthogonal.  In particular, every nontrivial coupled class
has a Killing leg on at least one factor, and no class with both legs
orthogonal to the Killing fields survives.
\end{theorem}

\begin{proof}
Both $\mathfrak G_h$ and $\mathfrak P_h$ lie in $\mathfrak K_h$.
For the first space this follows either from diffeomorphism invariance
or by linearizing the product curvature tensor on a mixed four-tuple;
for the second it follows by varying the two factor metrics separately.

Write an arbitrary symmetric tensor in block form as
\[
 B=P+B_C+Q,
\]
where $P$ and $Q$ are families of pure symmetric tensors on the first
and second factors.  Apply the usual divergence--Killing-operator
decomposition on each factor, with the other variable as a parameter.
After subtracting a Lie derivative, we may suppose
\begin{equation}\label{eq:partial-divergence-gauge}
 \delta_1^{(2)}P=0,
 \qquad
 \delta_2^{(2)}Q=0,
\end{equation}
where $\delta_i^{(2)}=\mathcal D_i^*$ is the divergence on symmetric
$2$-tensors.  The elliptic solution can be chosen orthogonal to
$\mathcal K_i$, so it depends smoothly on the parameter from the
other factor.

Polarizing $\mathfrak m_B(X,Y)$ separately in $X$ and $Y$ gives a
section $\mathcal M(B)$ of
$S^2T^*M_1\boxtimes S^2T^*M_2$.  In the block notation above,
\begin{equation}\label{eq:mixed-operator-block-form}
 \mathcal M(B)
 =\mathcal D_1\mathcal D_2C
  -\frac12\Hess_1Q-\frac12\Hess_2P.
\end{equation}
Here, for example,
$(\Hess_1Q)(X,X;Y,Y)=(\nabla^2_{X,X}Q)(Y,Y)$.
In adapted indices this is the pointwise identity
\[
 \mathcal M(B)_{ij\alpha\beta}
 =\nabla_{(i}\nabla_{(\alpha}C_{j)\beta)}
  -\frac12\nabla_i\nabla_jQ_{\alpha\beta}
  -\frac12\nabla_\alpha\nabla_\beta P_{ij},
\]
where the two pairs of parentheses are symmetrized independently.
Because product derivatives commute,
\eqref{eq:partial-divergence-gauge} implies that $\Hess_1Q$ is in
$\ker\delta_2^{(2)}$ in its second-factor slots.  The other two terms
in~\eqref{eq:mixed-operator-block-form} are in
$\operatorname{im}\mathcal D_2$: the last one is the second-factor
Hessian of the scalar $P(X,X)$.  Orthogonality of
$\operatorname{im}\mathcal D_2$ and $\ker\delta_2^{(2)}$ therefore
gives $\Hess_1Q=0$.  Interchanging the factors then gives
$\Hess_2P=0$, and hence
\begin{equation}\label{eq:pure-mixed-killing-equation}
 \mathcal D_1\mathcal D_2C=0.
\end{equation}
On a closed connected manifold a function with zero Hessian is
constant.  Applied to the scalar coefficients of $P$ and $Q$ in the
opposite factor, this shows that $P$ and $Q$ are pulled back from the
individual factors.  They may therefore be absorbed into
$\mathfrak P_h$.

It remains to solve~\eqref{eq:pure-mixed-killing-equation}.  Let
$L_i=\mathcal D_i^*\mathcal D_i$.  This is a nonnegative
self-adjoint elliptic operator on one-forms, with kernel
$\mathcal K_i$.  Choose smooth orthonormal eigenbases
$\{\alpha_r\}$ and $\{\beta_s\}$, with eigenvalues $\lambda_r$ and
$\mu_s$, and expand $C=\sum c_{rs}\alpha_r\otimes\beta_s$.
Then
\begin{equation}\label{eq:killing-spectral-product}
 \norm{\mathcal D_1\mathcal D_2C}_{L^2}^2
 =\sum_{r,s}\lambda_r\mu_s|c_{rs}|^2.
\end{equation}
Consequently $c_{rs}=0$ whenever both eigenvalues are positive, and
\begin{equation}\label{eq:raw-coupled-kernel}
 C\in(\mathcal K_1\otimes\Omega^1(M_2))
       +(\Omega^1(M_1)\otimes\mathcal K_2).
\end{equation}
The smooth spectral projections onto the finite-dimensional Killing
spaces show that the sums in~\eqref{eq:raw-coupled-kernel} may be
taken to be finite in the Killing index.

Writing $\Pi_i$ for the $L^2$ projection onto $\mathcal K_i$, choose
the nonredundant decomposition
$C=(\Pi_1\otimes I)C+((I-\Pi_1)\otimes\Pi_2)C$.
Use the Hodge splittings
$\Omega^1(M_i)=dC^\infty(M_i)\oplus\mathcal Z_i$ on the non-Killing
coefficients.  If $\kappa\in\mathcal K_1$,
$\lambda\in\mathcal K_2$, and $u,v$ are smooth functions, then
\begin{equation}\label{eq:coupled-exact-gauge}
 B_{\kappa\otimes du}=\mathcal L_{u\kappa^\sharp}h,
 \qquad
 B_{dv\otimes\lambda}=\mathcal L_{v\lambda^\sharp}h.
\end{equation}
Thus every class has a representative in~\eqref{eq:canonical-coupled-space}.

Such a representative is trace-free and divergence-free: in each
summand one leg is Killing and the other is co-closed.  Hence $B_C$
is $L^2$-orthogonal to every Lie derivative.  It is also pointwise
orthogonal to every pure product tensor.  Therefore no nonzero
$B_C$ in $\mathfrak H_{12}$ belongs to
$\mathfrak G_h+\mathfrak P_h$, proving uniqueness and
\eqref{eq:full-kernel-quotient}.  Finally,
$(\mathcal K_1\otimes\mathcal Z_2)\cap
 (\mathcal Z_1\otimes\mathcal K_2)
=\mathcal K_1\otimes\mathcal K_2$, which gives the orthogonal
splitting~\eqref{eq:coupled-orthogonal-splitting}.
\end{proof}

\begin{corollary}[The sphere-product kernel]
\label{cor:sphere-coupled-kernel}
Let $\mathfrak k_m\subset\Omega^1(\Sph^m)$ be the metric duals of
the rotational fields $x\mapsto Ax$, $A\in\mathfrak{so}(m+1)$, and
put
\[
 \mathcal Z_m^{\geq2}=\ker d^*\ominus\mathfrak k_m.
\]
Then the coupled two-jet kernel on $\Sph^m\times\Sph^n$, modulo
diffeomorphisms and factor tensors, is exactly
\begin{equation}\label{eq:sphere-coupled-kernel}
 (\mathfrak k_m\otimes\mathfrak k_n)
 \oplus(\mathfrak k_m\otimes\mathcal Z_n^{\geq2})
 \oplus(\mathcal Z_m^{\geq2}\otimes\mathfrak k_n).
\end{equation}
For $\Sph^2\times\Sph^2$, in coexact vector spherical harmonics,
the allowed bidegrees are precisely
\begin{equation}\label{eq:s2-harmonic-bidegrees}
 (1,1),\qquad (1,\ell)\ (\ell\geq2),
 \qquad (\ell,1)\ (\ell\geq2).
\end{equation}
With the cutoff $\ell\leq L$ on both factors, the quotient therefore
has dimension
\begin{equation}\label{eq:s2-kernel-cutoff-dimension}
 9+6\sum_{\ell=2}^{L}(2\ell+1)
 =6(L+1)^2-15.
\end{equation}
\end{corollary}

\begin{proof}
Every Killing field on a round sphere is rotational, and
$\dim\mathfrak k_m=\binom{m+1}{2}$.  Thus
Theorem~\ref{thm:full-coupled-kernel} gives
\eqref{eq:sphere-coupled-kernel}.  On $\Sph^2$, the Hodge star writes
every co-closed one-form as $*du$ modulo a constant potential.  Expanding
$u$ in scalar spherical harmonics, degree $1$ gives exactly the three
rotational Killing forms, while degree $\ell$ has multiplicity
$2\ell+1$.  This proves~\eqref{eq:s2-harmonic-bidegrees}, and
$\sum_{\ell=2}^{L}(2\ell+1)=(L+1)^2-4$, proving the dimension formula.
\end{proof}

Thus the full quadratic search space is infinite-dimensional, but it
is spectrally thin: one of the two harmonic legs must be rotational.
The rotational tensor~\eqref{eq:rotational-kernel} below lies in the
finite $9$-dimensional $(1,1)$ block.  The next theorem tests the two
one-sided towers $(1,\ell)$ and $(\ell,1)$, $\ell\geq2$, level by
level and without a separability assumption.

\subsection{Exact nonlinear tests in the coupled kernel}

Nearby moving planes already appear in the classical higher-order analysis
of Bourguignon--Deschamps--Sentenac
\cite{BourguignonDeschampsSentenac1973}.  For the round sphere, the
Killing-operator normal form and the spherical spectral gap make their sign
explicit on every one-sided harmonic level.

Identify the Killing one-forms on the first $\Sph^2$ with $\R^3$ by
$a\mapsto\kappa_a$, where $\kappa_a$ is dual to $x\mapsto a\times x$.
For an $\R^3$-valued one-form $E$ on the second sphere, define
\begin{equation}\label{eq:one-sided-tensor}
 C_E(X,Y)=\ip{x\times X}{E_y(Y)}.
\end{equation}
Thus $B_{C_E}$ is an arbitrary element of
$\mathfrak k_2\otimes\Omega^1(\Sph^2)$ in the notation of
Corollary~\ref{cor:sphere-coupled-kernel}.

\begin{lemma}[The moving-plane coefficient]
\label{lem:one-sided-moving-plane}
Fix $y\in\Sph^2$, choose a unit vector
$x\perp\operatorname{im}E_y$, and let $v\in x^\perp$ be unit.  Choose
unit $X\in T_x\Sph^2$ so that $x\times X=v$.  For an oriented
orthonormal pair $Y,Z\in T_y\Sph^2$, set
\[
 \theta=\ip{v}{E},\qquad f=*d\theta,
\]
where the star is on the second sphere.  For $c\in\R$, let
\begin{equation}\label{eq:one-sided-moving-plane}
 \Pi_s(c)=\spanop\{(X,scZ),(0,Y)\}.
\end{equation}
Then, as $s\to0$,
\begin{equation}\label{eq:one-sided-moving-coefficient}
 \secop_{h+sB_{C_E}}\Pi_s(c)
 =s^2\left(c^2-c(Yf)+\frac14f^2\right)+O(s^3).
\end{equation}
In particular, the choice $c=(Yf)/2$ gives
\begin{equation}\label{eq:one-sided-optimized-coefficient}
 \secop_{h+sB_{C_E}}\Pi_s\left(\frac{Yf}{2}\right)
 =\frac{s^2}{4}\bigl(f^2-(Yf)^2\bigr)+O(s^3).
\end{equation}
\end{lemma}

\begin{proof}
Use oriented product-normal coordinates at $(x,y)$ with coordinate
vectors $X$ on the first factor and $Y,Z$ on the second.  The condition
$x\perp\operatorname{im}E_y$ says that the first-factor curl of the
combined Killing one-form $C_E(\,\cdot\,,Y)$ vanishes at the base
point.  The fixed mixed-plane term is therefore
$s^2f^2/4$: it is the squared Christoffel term in
\eqref{eq:separated-exact-mixed}, with the same calculation applied
linearly to the sum of Killing legs.  The round curvature of the
second factor contributes $s^2c^2$ when $scZ$ is added to the first
spanning vector.  Finally, the term bilinear in the metric variation
and the plane displacement is $-s^2c(Yf)$.  Equivalently, at the
normal-coordinate origin write
$\theta=\theta_1\,dy^1+\theta_2\,dy^2$.  The three terms are
\[
 \frac{s^2}{4}(\partial_1\theta_2-\partial_2\theta_1)^2,
 \qquad s^2c^2,
 \qquad
 -s^2c\,\partial_1
   (\partial_1\theta_2-\partial_2\theta_1).
\]
The squared area is $1+O(s^2)$, proving
\eqref{eq:one-sided-moving-coefficient}.  Completing the square proves
\eqref{eq:one-sided-optimized-coefficient}.
\end{proof}

Let $\mathcal Z_\ell\subset\Omega^1(\Sph^2)$ denote the coexact
one-form spherical harmonics of degree $\ell$.

\begin{theorem}[Every fixed one-sided level has a quadratic oblique obstruction]
\label{thm:one-sided-towers}
For every $\ell\geq2$ and every nonzero
\[
 C\in\mathfrak k_2\otimes\mathcal Z_\ell,
\]
there are a point, a mixed plane at that point, a family of nearby planes
$\Pi_s$, and a constant $\delta_C>0$ such that
\begin{equation}\label{eq:one-sided-negative}
 \secop_{h+sB_C}(\Pi_s)=-\delta_Cs^2+O(s^3).
\end{equation}
Consequently $h+sB_C$ has negative sectional curvature for every
sufficiently small nonzero $s$.  The same conclusion holds for every
nonzero $C\in\mathcal Z_\ell\otimes\mathfrak k_2$.
\end{theorem}

\begin{proof}
Write $C=C_E$ as in~\eqref{eq:one-sided-tensor}, with each component
of $E$ in $\mathcal Z_\ell$, and put
\[
 F=*dE:\Sph^2\longrightarrow\R^3.
\]
The function $F$ is a nonzero vector-valued spherical harmonic of
degree $\ell$.  With $\lambda_\ell=\ell(\ell+1)$,
\begin{equation}\label{eq:vector-harmonic-energy}
 \int_{\Sph^2}|dF|^2
 =\lambda_\ell\int_{\Sph^2}|F|^2.
\end{equation}
If $|dF_y|_{\mathrm{op}}\leq|F(y)|$ at every point, then
$|dF|^2\leq2|F|^2$ pointwise, contradicting
\eqref{eq:vector-harmonic-energy} because
$\lambda_\ell\geq6>2$.  Hence there are $y$ and a unit $Y$ such that
\begin{equation}\label{eq:vector-harmonic-gap-point}
 |dF_y(Y)|>|F(y)|.
\end{equation}

Set $v=dF_y(Y)/|dF_y(Y)|$.  Componentwise Hodge theory gives
$E=\pm\lambda_\ell^{-1}*dF$, so
$\operatorname{im}E_y=\operatorname{im}dF_y$.  Choose a unit
$x\perp\operatorname{im}E_y$ and $X$ with $x\times X=v$.  For
$\theta=\ip{v}{E}$ and $f=*d\theta$, one has
\[
 f(y)=\ip{v}{F(y)},
 \qquad
 Yf=\ip{v}{dF_y(Y)}=|dF_y(Y)|.
\]
Thus~\eqref{eq:vector-harmonic-gap-point} implies $|Yf|>|f|$.
Lemma~\ref{lem:one-sided-moving-plane}, with $c=(Yf)/2$, gives
\eqref{eq:one-sided-negative} with
\[
 \delta_C=\frac14\bigl((Yf)^2-f^2\bigr)>0.
\]
Interchanging the two sphere factors proves the final assertion.
\end{proof}

\begin{lemma}[The two-curl moving-plane coefficient]
\label{lem:two-curl-moving-plane}
Let $C$ represent an element of the coupled mixed two-jet kernel on
$\Sph^2\times\Sph^2$.  Fix oriented orthonormal ambient frames
\[
 (x,X,\bar X),\qquad (y,Y,\bar Y),
 \qquad
 \bar X=x\times X,\quad \bar Y=y\times Y.
\]
For the two scalar one-forms obtained by freezing one argument of $C$,
put
\begin{equation}\label{eq:two-factor-curls}
 \alpha=*_{1}d_1\bigl(C(\,\cdot\,,Y)\bigr),
 \qquad
 \beta=*_{2}d_2\bigl(C(X,\,\cdot\,)\bigr).
\end{equation}
For $a,b\in\R$, let
\begin{equation}\label{eq:two-parameter-moving-plane}
 \Pi_s(a,b)=
 \spanop\{(X,sa\bar Y),(sb\bar X,Y)\}.
\end{equation}
Then
\begin{equation}\label{eq:two-curl-coefficient}
 \secop_{h+sB_C}\Pi_s(a,b)
 =s^2\left(
 a^2+b^2-a(Y\beta)-b(X\alpha)
 +\frac14(\alpha^2+\beta^2)\right)+O(s^3).
\end{equation}
In particular, $a=(Y\beta)/2$ and $b=(X\alpha)/2$ give
\begin{equation}\label{eq:two-curl-optimized}
 \secop_{h+sB_C}\Pi_s
 =-\frac{s^2}{4}\mathcal D_C+O(s^3),
 \qquad
 \mathcal D_C:=(X\alpha)^2+(Y\beta)^2-\alpha^2-\beta^2.
\end{equation}
\end{lemma}

\begin{proof}
Use product-normal coordinates adapted to the two displayed frames.
Since $C$ is in the mixed two-jet kernel, the
term linear in $s$ on the fixed mixed plane vanishes.  The round
curvatures of the two factors contribute $s^2a^2$ and $s^2b^2$.
The terms bilinear in the metric variation and the two plane
displacements are respectively
\[
 -s^2a(Y\beta),\qquad -s^2b(X\alpha).
\]
Finally, the first variation of the Christoffel symbol on the fixed
mixed plane has orthogonal factor components $\alpha/2$ and
$\beta/2$.  Its squared norm contributes
$s^2(\alpha^2+\beta^2)/4$.  The squared area is $1+O(s)$, so it does
not change the quadratic coefficient.  This proves
\eqref{eq:two-curl-coefficient}; minimizing its two independent
squares proves~\eqref{eq:two-curl-optimized}.
\end{proof}

\begin{theorem}[The full coupled kernel has a quadratic oblique obstruction]
\label{thm:full-kernel-obstruction}
Let $0\neq C\in\mathfrak H_{12}$ on
$\Sph^2\times\Sph^2$, where $\mathfrak H_{12}$ is the full coupled
kernel in~\eqref{eq:canonical-coupled-space}.  Then there are a point,
a background-mixed plane at that point, a family of nearby planes
$\Pi_s$, and a constant $\delta_C>0$ such that
\begin{equation}\label{eq:full-kernel-negative}
 \secop_{h+sB_C}(\Pi_s)=-\delta_Cs^2+O(s^3).
\end{equation}
Consequently $h+sB_C$ has negative sectional curvature for every
sufficiently small nonzero $s$.
\end{theorem}

\begin{proof}
The orthogonal harmonic splitting in
\eqref{eq:coupled-orthogonal-splitting} gives a unique decomposition
\begin{equation}\label{eq:full-harmonic-decomposition}
 C=C_M+C_E+C_D,
\end{equation}
where
\begin{align*}
 C_M(X,Y)&=\ip{x\times X}{M(y\times Y)},\\
 C_E(X,Y)&=\ip{x\times X}{E_y(Y)},\\
 C_D(X,Y)&=\ip{D_x(X)}{y\times Y}.
\end{align*}
Here $M:\R^3\to\R^3$ is linear, while the components of the
$\R^3$-valued coexact one-forms $E$ and $D$ are $L^2$-orthogonal to
the degree-one harmonics.  Put
\[
 F=*dE:\Sph^2\to\R^3,
 \qquad
 G=*dD:\Sph^2\to\R^3.
\]

Average the defect $\mathcal D_C$ in
\eqref{eq:two-curl-optimized} over the two oriented orthonormal frame
bundles, using probability Haar measure, and write
$\operatorname{Av}_{\Sph^2}$ for normalized area average.  The
isotropic frame identities
\[
 \operatorname{Av}\ip{x}{z}^2
 =\operatorname{Av}\ip{X}{z}^2
 =\operatorname{Av}\ip{\bar X}{z}^2
 =\frac13|z|^2
\]
and their second-factor analogues give
\begin{align}
 \operatorname{Av}\mathcal D_C
 =\frac13\bigg\{&
 \frac12\operatorname{Av}_{\Sph^2}|dF|^2
 -\operatorname{Av}_{\Sph^2}|F|^2 \notag\\
 &+\frac12\operatorname{Av}_{\Sph^2}|dG|^2
 -\operatorname{Av}_{\Sph^2}|G|^2
 \bigg\}.                                      \label{eq:full-defect-average}
\end{align}
For completeness, the two factorwise curls used here are
\begin{align*}
 \alpha={}&2\ip{x}{M\bar Y}
            +2\ip{x}{E(Y)}+\ip{G(x)}{\bar Y},\\
 \beta={}&2\ip{\bar X}{My}
            +\ip{\bar X}{F(y)}+2\ip{y}{D(X)},
\end{align*}
and their relevant derivatives are
\begin{align*}
 X\alpha={}&2\ip{X}{M\bar Y}
              +2\ip{X}{E(Y)}+\ip{dG_x(X)}{\bar Y},\\
 Y\beta={}&2\ip{\bar X}{MY}
              +\ip{\bar X}{dF_y(Y)}+2\ip{Y}{D(X)}.
\end{align*}
These formulas also verify~\eqref{eq:full-defect-average} directly.
The $M$-square has zero average.  Cross terms between different
summands vanish: the higher coexact one-forms are orthogonal to the
Killing one-forms, and their curls contain no degree-one scalar
harmonics.  The remaining pure $E$ and $D$ terms are exactly the two
lines on the right-hand side of~\eqref{eq:full-defect-average}.

Every component of $F$ and $G$ is a sum of scalar spherical harmonics
of degrees at least two.  Hence the spherical spectral gap gives
\[
 \int|dF|^2\geq6\int|F|^2,
 \qquad
 \int|dG|^2\geq6\int|G|^2.
\]
Thus~\eqref{eq:full-defect-average} is strictly positive whenever
$E\neq0$ or $D\neq0$.  Some pair of frames then has
$\mathcal D_C>0$, and Lemma~\ref{lem:two-curl-moving-plane} gives
\eqref{eq:full-kernel-negative}.

It remains to treat $E=D=0$ and $M\neq0$.  Choose oriented signed
singular frames with $Mv_i=\tau_i u_i$ and $\tau_1\neq0$, and set
\[
 (x,X,\bar X)=(u_2,u_1,-u_3),
 \qquad
 (y,Y,\bar Y)=(v_2,v_3,v_1).
\]
Then $\alpha=\beta=0$, while
\[
 X\alpha=2\tau_1,
 \qquad
 Y\beta=-2\tau_3.
\]
Consequently
$\mathcal D_{C_M}=4(\tau_1^2+\tau_3^2)>0$, and the same lemma finishes
the proof.
\end{proof}

\begin{corollary}[Every coupled kernel ray fails quartically]
\label{cor:coupled-kernel-repair}
Let $0\neq C\in\mathfrak H_{12}$, let
$\varphi\in C^\infty(\Sph^2\times\Sph^2)$, put $W=\nabla\varphi$,
and let $\Psi_t$ be the flow of $W/2$.  Then
\begin{equation}\label{eq:coupled-kernel-corrected-path}
 g_t=\Psi_t^*(h+t^2B_C)
 =h+t\Hess\varphi+t^2
   \left(\frac18\mathcal L_W^2h+B_C\right)+O(t^3)
\end{equation}
passes every quadratic mixed repair identity, but contains planes with
\begin{equation}\label{eq:coupled-kernel-quartic-obstruction}
 \secop_{g_t}=-\delta_Ct^4+O(t^6)<0
\end{equation}
for all sufficiently small nonzero $t$.
\end{corollary}

\begin{proof}
The expansion is the same pullback calculation as in
\eqref{eq:gauge-repair}.  The normalized second-order tensor is $B_C$,
which belongs to the mixed two-jet kernel.  Pullback preserves sectional
curvature; substitute $s=t^2$ in~\eqref{eq:full-kernel-negative}.
\end{proof}

Theorem~\ref{thm:full-kernel-obstruction} includes cross-degree
superpositions, simultaneous mixtures of both one-sided towers, and
arbitrary additions from the full $(1,1)$ block.  For comparison,
separated Killing variations already
occur in the classical variation theory
\cite{BourguignonDeschampsSentenac1973}.  We now record their exact
fixed-mixed curvature formula and then pass to the canonical transverse
sum in the $(1,1)$ block.  For one-forms $\xi$ on $M_1$ and $\eta$ on
$M_2$, define
the pure mixed symmetric tensor $B_{\xi,\eta}$ by
\begin{equation}\label{eq:separated-kernel-tensor}
 B_{\xi,\eta}((X,Y),(X',Y'))
 =\xi(X)\eta(Y')+\xi(X')\eta(Y).
\end{equation}

\begin{proposition}[Separated coupled kernel]
\label{prop:separated-kernel}
For every background-mixed unit pair $X,Y$,
\begin{equation}\label{eq:separated-kernel-factorization}
 \mathfrak m_{B_{\xi,\eta}}(X,Y)
 =(\nabla_X\xi)(X)(\nabla_Y\eta)(Y).
\end{equation}
Consequently $B_{\xi,\eta}$ lies in the mixed two-jet kernel if and
only if at least one of $\xi$ and $\eta$ is a Killing one-form.
If $\xi$ is Killing and $\eta=du$, then
\begin{equation}\label{eq:exact-separated-gauge}
 B_{\xi,du}=\mathcal L_{u\xi^\sharp}h,
\end{equation}
so the exact part of the second one-form is gauge.

Suppose $\xi$ is Killing and $h+sB_{\xi,\eta}$ is Riemannian.  Put
\begin{equation}\label{eq:separated-CXY}
 \mathcal C_{XY}
 =\frac12\bigl(\eta(Y)d\xi(X,\cdot),
                 \xi(X)d\eta(Y,\cdot)\bigr).
\end{equation}
Then every background-mixed plane satisfies the exact identity
\begin{equation}\label{eq:separated-exact-mixed}
 \secop_{h+sB_{\xi,\eta}}(X\wedge Y)
 =\frac{s^2}{\Delta_s(X,Y)}
  (h+sB_{\xi,\eta})^{-1}(\mathcal C_{XY},\mathcal C_{XY})
 \geq0.
\end{equation}
If $M_1=\Sph^2$, $\xi\neq0$, and $\dim M_2\geq2$, every such ray
nevertheless has a background-mixed zero-curvature plane.
\end{proposition}

\begin{proof}
Use product-normal extensions of $X$ and $Y$.  The two pure blocks of
$B_{\xi,\eta}$ vanish, while differentiating its mixed block twice
gives~\eqref{eq:separated-kernel-factorization}.  A one-form $\xi$ is
Killing exactly when $(\nabla_X\xi)(X)=0$ for every $X$.  Since the
two factors in~\eqref{eq:separated-kernel-factorization} depend on
independent variables, their product vanishes identically exactly when
one of the one-forms is Killing.  Formula
\eqref{eq:exact-separated-gauge} follows directly from the Killing
equation.

Let $\mathcal C$ be the first variation of the Christoffel symbol of
the first kind.  The Killing equation gives
\[
 \mathcal C_{XX}=0,
 \qquad
 \mathcal C_{XY}
 =\frac12\bigl(\eta(Y)d\xi(X,\cdot),
                 \xi(X)d\eta(Y,\cdot)\bigr).
\]
The second-derivative part of the lower-index mixed curvature is
linear in $s$ and vanishes by
\eqref{eq:separated-kernel-factorization}.  The exact remaining term
is therefore
\[
 s^2(h+sB_{\xi,\eta})^{-1}(\mathcal C_{XY},\mathcal C_{XY}),
\]
because the term pairing $\mathcal C_{XX}$ with
$\mathcal C_{YY}$ is zero.  Division by the squared area proves
\eqref{eq:separated-exact-mixed}.

A nonzero Killing field on the round $2$-sphere is rotational and has
a zero $x$.  At $(x,y)$ choose a nonzero $Y\in\ker\eta_y$, which is
possible when $\dim M_2\geq2$.  Then~\eqref{eq:separated-CXY} vanishes
for every $X$, giving the final assertion.
\end{proof}

The exact parts in~\eqref{eq:exact-separated-gauge} disappear in the
quotient by diffeomorphisms.  On a sphere the Hodge decomposition
therefore leaves the co-closed part of the other one-form as the first
genuinely coupled separated sector.  The following finite sum is the
most symmetric such tensor on $\Sph^2\times\Sph^2$.

Identify the two ambient Euclidean spaces with $\R^3$.  For
$a\in\R^3$, let $K_a(x)=a\times x$ and $L_a(y)=a\times y$, and denote
their metric-dual one-forms by $\kappa_a$ and $\lambda_a$.  For an
orthonormal basis $a_0,a_1,a_2$ of $\R^3$, set
\begin{equation}\label{eq:rotational-kernel}
 B_{\mathrm{rot}}=\sum_{r=0}^2B_{\kappa_{a_r},\lambda_{a_r}}.
\end{equation}
Equivalently, its only nonzero block is
\begin{equation}\label{eq:rotational-kernel-intrinsic}
 B_{\mathrm{rot}}((X,0),(0,Y))
 =\ip{x\times X}{y\times Y}.
\end{equation}

\begin{theorem}[A nongauge kernel ray with an oblique obstruction]
\label{thm:rotational-kernel}
The tensor $B_{\mathrm{rot}}$ is trace-free and divergence-free.  It is
nonzero and $L^2$-orthogonal to all Lie derivatives and to all tensors
pulled back from the individual factors.  In particular, it represents
a genuinely coupled, nongauge element of the mixed two-jet kernel.

The metric
\[
 k_s=h+sB_{\mathrm{rot}}
\]
has the sharp Riemannian range
\[
 k_s>0\quad\Longleftrightarrow\quad |s|<1.
\]
Throughout this interval every
background-mixed sectional curvature is nonnegative.  At each diagonal
point $(x,x)$ every background-mixed sectional curvature is zero.
Nevertheless, fix an orthonormal basis $(e_0,e_1,e_2)$ of $\R^3$ and
put $p=(e_0,e_0)$.  For $i=1,2$, define
\[
 E_i^+=\frac{(e_i,e_i)}{\sqrt{2(1+s)}},
 \qquad
 E_i^-=\frac{(e_i,-e_i)}{\sqrt{2(1-s)}}.
\]
The $k_s$-orthonormal plane spanned by
\begin{equation}\label{eq:rotational-oblique-plane}
 U=\frac{E_1^+-E_1^-}{\sqrt2},
 \qquad
 V=\frac{E_2^++E_2^-}{\sqrt2}
\end{equation}
has
\begin{equation}\label{eq:rotational-oblique-curvature}
 \secop_{k_s}(U\wedge V)
 =-\frac{3s^2}{2(1-s)^2(1+s)}<0
 \qquad(0<|s|<1).
\end{equation}
Thus this first transverse coupled kernel ray leaves the
nonnegative-curvature cone immediately, despite improving no fixed
mixed plane in the negative direction.
\end{theorem}

\begin{proof}
Each $\kappa_{a_r}$ and $\lambda_{a_r}$ is Killing, so
Proposition~\ref{prop:separated-kernel} and linearity give
$\mathfrak m_{B_{\mathrm{rot}}}=0$.  The tensor is pure mixed, hence
trace-free.  Its divergence vanishes because every Killing one-form is
co-closed.  Integration by parts makes a divergence-free symmetric
tensor $L^2$-orthogonal to every Lie derivative; pointwise block
orthogonality gives the assertion for product tensors.

The maps $X\mapsto x\times X$ and $Y\mapsto y\times Y$ are isometries
on their tangent planes.  Hence
\[
 |B_{\mathrm{rot}}((X,0),(0,Y))|
 \leq\norm{X}\norm{Y}.
\]
It follows that $k_s>0$ for $|s|<1$.  Equality is attained in this
estimate when $x=y$ and $X=Y$, so a symmetric or antisymmetric vector
becomes null at $s=-1$ or $s=1$, respectively; outside the closed
interval the metric is indefinite.  This proves the sharp range.

The proof of~\eqref{eq:separated-exact-mixed} applies unchanged to a
sum whose first one-forms are Killing, proving mixed nonnegativity.
For a rotational Killing form,
\[
 d\kappa_a(X,Z)=2\ip{a}{X\times Z}.
\]
At a diagonal point the two components of the summed covector
$\mathcal C_{XY}$ contain scalar products between a tangent vector of
the form $x\times Y$ and a normal vector of the form $X\times Z$;
they vanish.  Hence all background-mixed curvatures there are zero.

Finally use orthographic coordinates
\[
 x=(\sqrt{1-|u|^2},u_1,u_2),
 \qquad
 y=(\sqrt{1-|v|^2},v_1,v_2)
\]
centered at $p$.  The mixed block is
\[
 (B_{\mathrm{rot}})_{i\alpha}
 =\ip{x\times\partial_{u_i}x}
       {y\times\partial_{v_\alpha}y}.
\]
Substitution into the lower-index curvature formula for the vectors
in~\eqref{eq:rotational-oblique-plane} gives unit squared area and
\eqref{eq:rotational-oblique-curvature}.
\end{proof}

\begin{corollary}[The first nongauge repair fails quartically]
\label{cor:rotational-repair}
Let $f\in C^\infty(\Sph^2\times\Sph^2)$, put $Z=\nabla f$, and let
$\Psi_t$ be the flow of $Z/2$.  Then
\begin{equation}\label{eq:rotational-corrected-path}
 g_t=\Psi_t^*(h+t^2B_{\mathrm{rot}})
 =h+t\Hess f+t^2
  \left(\frac18\mathcal L_Z^2h+B_{\mathrm{rot}}\right)+O(t^3).
\end{equation}
This path satisfies every quadratic mixed repair identity in
Theorem~\ref{thm:integrated-two-jet}.  Nevertheless, for every
$0<|t|<1$ it contains a plane of curvature
\begin{equation}\label{eq:rotational-quartic-obstruction}
 -\frac{3t^4}{2(1-t^2)^2(1+t^2)}<0.
\end{equation}
\end{corollary}

\begin{proof}
The expansion follows from pullback by the flow.  Its normalized
second-order tensor is $B_{\mathrm{rot}}$, which lies in the mixed
two-jet kernel.  Pullback preserves sectional curvature, so substitute
$s=t^2$ in~\eqref{eq:rotational-oblique-curvature}.
\end{proof}

\section{Consequences, context, and scope}\label{sec:scope}

\subsection{The four-dimensional Hopf case}

On $\Sph^2\times\Sph^2$, every bilinear product of degree-one
spherical harmonics has the form $s_A$.  Orthogonal changes of
coordinates in the two factors reduce $A$ to singular-value form, so
Theorem~\ref{thm:main} is a classification-free result for the whole
nine-dimensional space of such potentials.  On the open dense stratum
$\sigma_1>\sigma_2$, Theorem~\ref{thm:spectral-gap} detects the
negative curvature on an explicit background-mixed plane.  On the
exceptional stratum $\sigma_1=\sigma_2$, the flat-torus theorem
supplies three embedded zero-curvature loci, while
Proposition~\ref{prop:top-critical-oblique} detects strict negative
curvature on an explicit oblique plane.
At the isotropic endpoint $\sigma_1=\sigma_2=\sigma_3$,
Theorem~\ref{thm:isotropic-mixed} goes further and proves that every
background-mixed curvature is nonnegative on the full Riemannian
interval.  Thus the isotropic endpoint makes the distinction between
background-mixed and oblique planes unavoidable.  Specializing
Corollary~\ref{cor:complete-isolation} shows that every nontrivial ray
in the whole nine-dimensional family leaves the
nonnegative-curvature cone immediately.  The top-critical oblique
mechanism works throughout rank at least two, the zero-critical
oblique mechanism handles rank one, and the simple-top stratum
additionally admits the background-mixed spectral-gap mechanism.
Corollary~\ref{cor:two-jet-cost} strengthens the local conclusion:
after the infinitesimal diffeomorphism has been removed, every
nonzero ray already has a negative quadratic mixed-curvature
coefficient.  Thus the quartic spectral-gap plane is a special
background-mixed phenomenon, not the first obstruction among all
planes.

This conclusion is independent of the Hsiang--Kleiner theorem.  In the
generic full-rank case the inherited background symmetry is finite, so
one cannot obtain the conclusion merely by pointing to an inherited
circle action; Theorem~\ref{thm:spectral-gap} in fact gives strict
negative curvature there.  Conversely,
Corollary~\ref{cor:curvature-symmetry} records that the repeated-top
case does retain an inherited circle action, but the explicit oblique
calculation already proves negative curvature there.  We emphasize
again that this does not prove the full isometry group is finite in the
generic case; it shows only that the direct curvature and flat-torus
arguments do not need continuous inherited symmetry there.

\subsection{Relation to deformations of product metrics}

Bourguignon's work placed deformations of product metrics into the
broader study of Hopf's conjecture
\cite{Bourguignon1975,BourguignonDeschampsSentenac1973}.  The family
considered here is a particular exact Hessian deformation through the
round product.  Its special feature is not a general stability theorem
for product metrics, but an algebraic compatibility between singular
two-planes and fixed-point geometry.  That compatibility persists for
the entire interval~\eqref{eq:pd-range}, including parameters far from
the infinitesimal regime.

The mixed-curvature calculation adds a complementary point.  On each
fixed background-mixed plane its second-order term is nonnegative,
which is compatible with the classical infinitesimal analysis of
product variations \cite{BourguignonDeschampsSentenac1973}.
Theorem~\ref{thm:spectral-gap} shows that when this coefficient
vanishes, the first negative term on a distinguished background-mixed
plane can be quartic.  Theorem~\ref{thm:two-jet} explains why this does
not make the full metric deformation second-order nonnegative: after
gauge normalization, a moving oblique plane becomes a fixed mixed
plane, and Corollary~\ref{cor:two-jet-cost} gives its negative
quadratic coefficient.  A nonlinear continuation can therefore
survive the two-jet test only if its correction tensor $B$ cancels the
quadratic mixed-curvature functional everywhere.  The integrated
identity~\eqref{eq:integrated-two-jet} forces equality, and the
singular planes give the explicit costs
\eqref{eq:two-jet-cost-rank-two} and
\eqref{eq:two-jet-cost-rank-one}.  The canonical gauge correction
\eqref{eq:gauge-repair} shows that this system is consistent, but by
itself produces no new geometry.
Theorem~\ref{thm:full-coupled-kernel} identifies the entire quadratic
kernel modulo gauge and product tensors; on $\Sph^2\times\Sph^2$ it
leaves only the bidegrees $(1,\ell)$ and $(\ell,1)$.
Theorem~\ref{thm:full-kernel-obstruction} tests every nonzero element
of that quotient at once: a moving plane has negative curvature at
order $s^2$, even for cross-degree superpositions, simultaneous
two-tower mixtures, and arbitrary $(1,1)$ components.  The
corresponding canonical corrected Hessian path fails at order $t^4$.
Proposition~\ref{prop:separated-kernel} gives the exact nonlinear
mixed-curvature formula for its separated elements.  The rotational
tensor in Theorem~\ref{thm:rotational-kernel} is transverse to the
diffeomorphism orbit and makes every fixed background-mixed curvature
nonnegative.  Its explicit oblique curvature is nevertheless negative
at order $s^2$, or order $t^4$ in the corrected Hessian path of
Corollary~\ref{cor:rotational-repair}.  Thus passing the complete
quadratic mixed test and the ensuing fixed-mixed quartic test is still
not enough to control nearby oblique planes.

Symmetry-based obstructions remain among the strongest general results
related to Hopf's question.  Hsiang and Kleiner proved that a positively
curved oriented four-manifold with an isometric circle action has
$b_2\leq1$ \cite{HsiangKleiner1989}; Amann and Kennard established
higher-dimensional product obstructions under quantitative torus
symmetry assumptions \cite{AmannKennard2017}.  Other authors have
treated special metric forms, including conformal and doubly warped
products \cite{SeckNiangMasiello2022}.  Our result is complementary:
the metrics generally have mixed Hessian blocks and need not be
conformal or warped, while their obstruction is an explicit flat
submanifold rather than a conclusion drawn from an effective torus
action.

\subsection{Higher-dimensional content and limitations}

The construction and the curvature isolation theorem are
dimension-independent.  For
$q=\min\{m+1,n+1\}$ it yields at least $\binom q2$ tori for each
singular frame, so the number of certified zero-curvature loci grows
quadratically with the smaller ambient dimension.  This gives a
family-level obstruction on every product of positive-dimensional
spheres.  When both spheres have dimension at least two, the
singular-subsphere reduction promotes the four-dimensional
top-critical calculation to every ambient dimension.  Combined with
the rank-one zero-critical calculation, it shows that every
nontrivial bilinear Hessian ray leaves the nonnegative-curvature cone
immediately, with the failure always detected inside the top-three singular
$\Sph^2\times\Sph^2$.  Corollary~\ref{cor:uniform-negative-margin}
upgrades this qualitative statement to an explicit dimension-free
negative bound on every normalized annulus away from the product
metric and the Riemannian boundary.

The result should not be read as a curvature classification of
$g_{\gamma,A}$.  We give an exact formula for every background-mixed
plane and determine a distinguished collection of zero sectional
curvatures.  Proposition~\ref{prop:top-critical-oblique} determines
one distinguished oblique plane for every matrix of rank at least two,
which is enough, together with the rank-one zero-critical plane, for
complete ray isolation in all dimensions at least two.  We do not
classify all planes whose spanning vectors both have components in
both factors.
Nor do Hessian deformations $h+\gamma\Hess s_A$ exhaust all metrics on
a product of spheres.  The two-jet theorem constrains arbitrary
nonlinear paths sharing the same Hessian first variation, and the
integrated theorem shows that a nonnegative path must cancel every
quadratic mixed-curvature coefficient.  This cancellation is feasible:
the gauge correction does it identically, and product variations lie
in the kernel as well.  Theorem~\ref{thm:full-coupled-kernel}
classifies the full quotient and shows that one harmonic leg must be
Killing.  Proposition~\ref{prop:separated-kernel} gives the exact
background-mixed curvature of separated representatives, and
Theorem~\ref{thm:full-kernel-obstruction} obstructs every nonzero
coupled kernel ray, including arbitrary mixtures of all allowed
harmonic sectors.
Theorem~\ref{thm:rotational-kernel} settles the canonical $(1,1)$ sum
by an exact formula.  What remains open is not a sector of the
quadratic kernel, but the higher-jet problem: arbitrary third- and
higher-order corrections can enter before or at the quartic
obstruction of a canonical corrected path.  A broader obstruction
would require compatibility conditions on those higher jets, a
substantially larger class of first variations, or a full curvature
analysis.  The present kernel-ray results are sharp in their stated
class.

\subsection{Verification}

The arguments combine the Hessian identity~\eqref{eq:hessian}, the
two-frame singular-value principle, the exact mixed-curvature
identity, and the fixed-point lemma for isometries.  As an independent
check against coordinate or sign errors,
the source distribution includes the script
\path{verify_hessian_tori.py}.  It constructs the induced torus
metric from~\eqref{eq:phi-uv}, computes its Levi-Civita connection and
Riemann tensor directly, and verifies symbolically that its Gaussian
curvature is zero.  It also checks the separated normal form and the
sharp endpoint degeneracies.  The companion script
\path{verify_mixed_curvature_counterexample.py} differentiates
the full orthographic-coordinate metric symbolically and checks the
mixed-curvature identity, the spectral-gap formula, the simple
zero-threshold specialization, the plane-wise nonzero threshold, and
the isotropic quaternion factorization, the uniform top-critical
oblique-plane curvature, its two-jet coefficient, and the rank-one
zero-critical specialization using exact arithmetic.  The separate
script \path{verify_killing_kernel.py} checks the diagonal rotational
kernel metric, its sharp Riemannian boundary, its flat diagonal mixed
plane, and its exact negative oblique curvature.  The script
\path{verify_coupled_kernel_classification.py} checks the factorization
through the two Killing operators, the tensor-product nullity formula,
and the gauge/co-closed normal form in an exact finite-mode model.  Finally,
\path{verify_coupled_kernel_obstruction.py} differentiates the exact
orthographic-coordinate metric, recovers the general two-parameter
moving-plane coefficient, tests a nonseparated degree-two vector
harmonic, and verifies the signed-SVD obstruction for an arbitrary
$(1,1)$ block.

\section*{Acknowledgments}

The author used large-language-model assistance during the process of
symbolic checking, editorial restructuring, and code preparation.  All
mathematical claims and proofs remain the responsibility of the author.

\begin{thebibliography}{10}

\bibitem{Amann2021}
M.~Amann,
\emph{Bounds on sectional curvature and interactions with topology},
Jahresber. Dtsch. Math.-Ver. \textbf{123} (2021), 27--55,
\href{https://doi.org/10.1365/s13291-020-00225-x}{doi:10.1365/s13291-020-00225-x}.

\bibitem{AmannKennard2017}
M.~Amann and L.~Kennard,
\emph{On a generalized conjecture of Hopf with symmetry},
Compos. Math. \textbf{153} (2017), no.~2, 313--322,
\href{https://doi.org/10.1112/S0010437X16008150}{doi:10.1112/S0010437X16008150}.

\bibitem{BourguignonDeschampsSentenac1972}
J.-P.~Bourguignon, A.~Deschamps, and P.~Sentenac,
\emph{Conjecture de H.~Hopf sur les produits de vari\'et\'es},
Ann. Sci. \'{E}cole Norm. Sup. (4) \textbf{5} (1972), 277--302,
\href{https://doi.org/10.24033/asens.1229}{doi:10.24033/asens.1229}.

\bibitem{BourguignonDeschampsSentenac1973}
J.-P.~Bourguignon, A.~Deschamps, and P.~Sentenac,
\emph{Quelques variations particuli\`eres d'un produit de m\'etriques},
Ann. Sci. \'{E}cole Norm. Sup. (4) \textbf{6} (1973), 1--16,
\href{https://doi.org/10.24033/asens.1240}{doi:10.24033/asens.1240}.

\bibitem{Bourguignon1975}
J.-P.~Bourguignon,
\emph{Some constructions related to H.~Hopf's conjecture on product manifolds},
in \emph{Differential Geometry}, Proc. Sympos. Pure Math., vol.~27,
part~1, Amer. Math. Soc., Providence, RI, 1975, pp.~33--37,
\href{https://doi.org/10.1090/pspum/027.1/9936}{doi:10.1090/pspum/027.1/9936}.

\bibitem{Bhatia1997}
R.~Bhatia,
\emph{Matrix Analysis}, Graduate Texts in Mathematics, vol.~169,
Springer, New York, 1997.

\bibitem{HsiangKleiner1989}
W.-Y.~Hsiang and B.~Kleiner,
\emph{On the topology of positively curved $4$-manifolds with symmetry},
J. Differential Geom. \textbf{29} (1989), no.~3, 615--621,
\href{https://doi.org/10.4310/jdg/1214443064}{doi:10.4310/jdg/1214443064}.

\bibitem{SeckNiangMasiello2022}
T.~Seck, A.~Niang, and A.~Masiello,
\emph{Conformal metrics to a product or doubly warped product on
$S^2\times S^2$ and the Hopf conjecture},
J. Math. \textbf{2022} (2022), Article ID 6268017,
\href{https://doi.org/10.1155/2022/6268017}{doi:10.1155/2022/6268017}.

\end{thebibliography}

\end{document}
